\documentclass[12pt,twoside]{ord}

%*******************AMS-latex*************
\usepackage{amssymb}
\usepackage{amsmath}     
\usepackage[T1]{fontenc}
\usepackage[utf8]{inputenc}
\usepackage{algorithmic}
\usepackage[]{algorithm2e}
\usepackage{lineno}
\usepackage{longtable}

\makeatletter
\providecommand{\FloatBarrier}{\par\ifx\@deferlist\@empty\else\clearpage\fi}
\makeatother

\emergencystretch 3em
\hfuzz=100pt
\vfuzz=100pt
\hbadness=10000
\vbadness=10000

%Assigned by authors
\ORDtitle{How Does Equity Return Predictability Vary Across Horizons and Market Regimes? An Out-of-Sample Analysis}
\ORDshortafill{D. D. H. Nguyen et al.}
\ORDshorttitle{How Does Equity Return Predictability Vary?}
\ORDaffilone{University of Information Technology, Vietnam National University Ho Chi Minh City, Ho Chi Minh City, Vietnam}
\ORDaffiltwo{}
\ORDaffilthree{}
\ORDaffilfour{}

\ORDauthorone{Doan Duc Hieu~Nguyen}{*1}{0009-0008-5897-8147}{24520500@gm.uit.edu.vn}
\ORDauthortwo{Thi Phuong Hoa~Nguyen}{1}{0009-0002-7378-5091}{24520524@gm.uit.edu.vn}
\ORDauthorthree{Van Phat~Nguyen}{1}{0009-0009-4400-0794}{24521307@gm.uit.edu.vn}
\ORDauthorfour{Anh Thu~Tran}{1}{0009-0000-7168-2123}{24521643@gm.uit.edu.vn}
\ORDauthorfive{Ngoc Thao~Tran}{1}{0009-0009-6563-9394}{24521609@gm.uit.edu.vn}
\ORDauthorsix{Phi Long~Duong}{1}{0000-0000-0000-0000}{longdp@uit.edu.vn}
\ORDcorrespoding{* author, email address:}{24520500@gm.uit.edu.vn}

%-------------------------------------------
\ORDnumber{x}
\ORDvolume{xx}
\ORDdoi{DOI: draft}
\ORDyear{2026}
\ORDReceived{xxxxx~xxxxx}
\ORDRevised{xxxxx~xxxxx}
\ORDAccepted{xxxxx~xxxxx}
\ORDPublished{xxxxx~xxxxx}
\ORDFirstpagenum{1}
\ORDFounding{}
%-------------------------------------------

\begin{document}

\maketitle

\begin{abstract}
This study investigates how out-of-sample (OOS) equity return predictability varies across forecast horizons, market volatility regimes, and chronological subperiods. Using daily data on the S\&P 500 index and the CBOE Volatility Index (VIX) spanning 1990--2026, we evaluate twelve representative forecasting models, covering autoregressive econometrics, penalized linear regressions, kernel-based methods, tree-based ensembles, and deep learning architectures, against the Historical Mean benchmark across five horizons (1D, 5D, 20D, 60D, and 120D). The forecasting framework incorporates VIX-derived predictors, while market conditions are classified into low, normal, and high volatility regimes based on the 20-day realized volatility of S\&P 500 returns. Three empirical findings emerge. First, forecasting performance follows a nonlinear horizon profile, with gains tending to concentrate at intermediate horizons (20D--60D), while becoming increasingly heterogeneous at 120D. Second, descriptive out-of-sample performance varies across volatility regimes and forecast horizons. Positive $R^2_{\mathrm{OOS}}$ is most prevalent under high volatility at the shortest horizon, while normal-volatility conditions exhibit broader predictive gains across multiple model classes, particularly at the 60D and 120D horizons. At 20D, however, positive results are more frequent under high volatility, indicating that regime-related performance patterns are not uniform across horizons. Formal interaction tests further show that statistically significant horizon--regime interactions are architecture-specific rather than widespread across models. Third, temporal stability varies across chronological subperiods, with sign reversals observed in selected model-horizon combinations. Overall, these findings suggest that equity return predictability is episodic and conditional on forecast horizon, market volatility, and historical environment, highlighting the limitations of static single-horizon evaluations.
\end{abstract}

\ORDKeywords{Equity return predictability, Multi-horizon forecasting, Volatility regimes, Machine learning, Out-of-sample evaluation}

\section{Introduction}\label{sec:intro}

The predictability of equity returns has long been a central question in financial economics, with important implications for asset pricing, portfolio allocation, and risk management. Despite extensive research, forecasting equity returns remains challenging because expected returns are small relative to return volatility, and predictive relationships may vary across investment horizons, market conditions, and historical periods~\cite{campbell1998,welch2008}. Although machine learning provides flexible approaches to capturing complex and potentially nonlinear relationships in financial data~\cite{gu2020}, greater model complexity does not necessarily translate into reliable out-of-sample (OOS) predictive performance. Consequently, understanding the conditions under which equity returns can be predicted remains as important as developing more sophisticated forecasting models.

Previous studies have documented several dimensions of equity return predictability. First, predictive performance may vary across forecast horizons. Nechv{\'a}talov{\'a}~\cite{nechvatalova2024} investigates multi-horizon equity return predictability using machine learning, while Xu~\cite{xu2020} examines alternative approaches to testing predictability across multiple horizons. These studies highlight the importance of considering the investment horizon when evaluating predictive relationships rather than relying on a single forecasting period. Second, market conditions, particularly volatility, may influence the effectiveness of return predictors. Baltas and Karyampas~\cite{baltas2018} emphasize the importance of regime-dependent evaluation in equity risk premium forecasting, while Cakici and Zaremba~\cite{cakici2025} show that the relative predictive value of accounting and technical information varies across investment horizons and market conditions. Related research also establishes the connection between market volatility, implied volatility, and expected stock returns~\cite{banerjee2007,bekaert2014,french1987}. Third, predictive relationships may change over time. Paye and Timmermann~\cite{paye2006} document substantial instability in return prediction models, while Harvey et al.~\cite{harvey2021} investigate the real-time detection of temporary periods of predictability in the US equity premium. Together, these findings suggest that equity return predictability should be evaluated across multiple horizons, market environments, and chronological periods.

Recent research has also highlighted the importance of rigorous OOS evaluation in financial forecasting. Campbell and Thompson~\cite{campbell2008} demonstrate the importance of comparing return forecasts against the historical average, while Welch and Goyal~\cite{welch2008} document the difficulty of obtaining reliable equity premium predictions across a broad range of predictive variables. More recently, Nu\~nez Perez et al.~\cite{perez2026} propose a leakage-controlled, horizon-specific model selection pipeline for daily equity forecasting. Their analysis considers short-term forecasting horizons and emphasizes the importance of chronological validation and controlling information leakage when evaluating predictive performance. These findings reinforce the need for careful evaluation designs, particularly when forecasting cumulative returns over multiple periods.

Despite these contributions, the existing literature examines different aspects of equity return predictability under varying empirical settings, data frequencies, forecasting objectives, and evaluation procedures. For example, although Cakici and Zaremba~\cite{cakici2025} investigate both forecast horizons and market conditions, their analysis focuses on the relative predictive value of accounting and technical information in stock return prediction. Nechv{\'a}talov{\'a}~\cite{nechvatalova2024} examines multi-horizon equity return predictability, whereas Baltas and Karyampas~\cite{baltas2018} focus on regime-dependent evaluation. Meanwhile, Nu\~nez Perez et al.~\cite{perez2026} provide a leakage-controlled framework for daily forecasting but concentrate on two short-term horizons. Paye and Timmermann~\cite{paye2006} and Harvey et al.~\cite{harvey2021} provide evidence of temporal instability and temporary predictive episodes, respectively, but address these issues through different empirical approaches. Therefore, further empirical evidence is useful on how predictive performance varies jointly across multiple forecast horizons and volatility regimes, and whether the resulting patterns remain consistent across chronological sub-periods within a common OOS evaluation framework for aggregate equity index returns.

To address this gap, this study investigates OOS equity return predictability using daily S\&P~500 index returns and CBOE Volatility Index (VIX) data spanning 1990--2026~\cite{fredvix,yahoosandp}. The empirical analysis evaluates 12 econometric and machine learning forecasting models across five forecast horizons: 1, 5, 20, 60, and 120 trading days. Predictive performance is examined under three market volatility regimes---low, normal, and high---and across chronological sub-periods to assess the temporal stability of the observed patterns. All forecasting models are evaluated against an expanding historical mean benchmark, following the importance of benchmark-based OOS assessment in equity return prediction~\cite{campbell2008}. The evaluation design also considers the importance of chronological forecasting procedures and appropriate treatment of multi-period forecasting targets~\cite{hewamalage2023,tashman2000}.

Specifically, this study addresses three research questions:
\begin{enumerate}
\item How does OOS equity return predictability vary across forecast horizons?
\item Does the relationship between forecast horizon and predictive performance differ across market volatility regimes?
\item Do the observed horizon- and regime-dependent patterns remain stable across chronological sub-periods?
\end{enumerate}

This study contributes to the empirical literature on equity return predictability in three ways. First, it provides a unified multi-model assessment of return predictability across five forecast horizons, comparing 12 forecasting models from different methodological families against a common historical mean benchmark. Second, it examines the heterogeneity of predictive performance across low-, normal-, and high-volatility regimes, providing evidence on whether the relationship between forecast horizon and predictive performance depends on prevailing market conditions. Third, it investigates the temporal stability of these horizon- and regime-dependent patterns across chronological sub-periods, allowing the analysis to distinguish predictive gains that are relatively persistent from those concentrated in particular periods. Rather than proposing a new forecasting algorithm, this study provides an integrated empirical evaluation of the horizon, regime, and temporal dimensions of equity return predictability within a common OOS framework.

The remainder of this paper is organized as follows. Section~\ref{sec:lit_review} reviews the related literature and develops the research hypotheses. Section~\ref{sec:data} describes the data and research design. Section~\ref{sec:methodology} presents the forecasting models and evaluation methodology. Section~\ref{sec:results} reports and discusses the empirical results, and Section~\ref{sec:conclusion} concludes the study.


\section{Related Work and Hypothesis Development}\label{sec:lit_review}

\subsection{Equity Return Predictability and Model Complexity}

Achieving reliable out-of-sample (OOS) equity return predictability beyond historical benchmarks remains a central challenge in empirical finance. Welch and Goyal~\cite{welch2008} document the limited and time-varying performance of conventional equity premium prediction models, highlighting the difficulty of identifying stable predictive relationships in aggregate stock returns. Campbell and Thompson~\cite{campbell2008} show that imposing economically motivated constraints can improve the OOS performance of predictive regressions. More recently, Gu et al.~\cite{gu2020} demonstrate that machine learning methods can capture nonlinear relationships and interactions in asset pricing, offering an alternative to conventional linear specifications.

However, greater model flexibility does not necessarily translate into superior OOS performance. When predictive signals are weak relative to return volatility, complex models may be more exposed to estimation uncertainty and overfitting. This creates a trade-off between model flexibility and generalization, motivating the comparison of models with different levels of complexity under a common OOS evaluation framework.

Beyond model specification, the evaluation of predictive performance must also account for the forecast horizon. Quaedvlieg~\cite{quaedvlieg2021} develops multi-horizon forecast comparison procedures that jointly evaluate predictive performance across a forecast path, rather than relying exclusively on separate single-horizon comparisons. This is particularly relevant because cumulative returns over overlapping horizons induce dependence between forecast errors. Xu~\cite{xu2020} further examines the statistical challenges of testing predictability across multiple horizons and shows that conventional inference procedures may become unreliable as the horizon increases. These considerations suggest that equity return predictability should be evaluated not only across competing model specifications but also across forecast horizons using procedures that account for horizon-related dependence.

\subsection{Horizon and Volatility-Regime Dependence}

The predictive performance of equity return models may vary across investment horizons because the relevance of financial information can differ between short- and long-term forecasting periods. Nechv{\'a}talov{\'a}~\cite{nechvatalova2024} documents substantial differences in predictive performance across monthly horizons in a machine learning framework for cross-sectional stock returns. Cakici and Zaremba~\cite{cakici2025} also investigate how the relative importance of accounting and technical information varies with the forecasting horizon. Although these studies provide evidence that predictive relationships are horizon-dependent, their settings differ from aggregate time-series forecasting of the equity premium. In particular, cross-sectional predictability does not necessarily imply the same horizon-related patterns in aggregate market returns.

The comparison of predictive performance across horizons also requires careful interpretation. A decline in statistical forecast accuracy at longer horizons does not necessarily imply a corresponding decline in economic value, since portfolio turnover and transaction costs may also change with the investment horizon. Nechv{\'a}talov{\'a}~\cite{nechvatalova2024}, for example, illustrates that statistical predictability and risk-adjusted portfolio performance need not exhibit identical patterns across horizons. Accordingly, the present study focuses on OOS predictive performance rather than assuming that longer horizons must produce uniformly weaker forecasting results.

Market volatility provides another important dimension of conditional predictability. Baltas and Karyampas~\cite{baltas2018} show that equity risk premium forecasting performance varies across market volatility regimes, with model performance differing substantially between low- and high-volatility conditions. Related evidence from Cakici and Zaremba~\cite{cakici2025} suggests that the predictive importance of information sources may also depend on volatility conditions. Other studies examine conditional predictability through alternative definitions of market states. For example, Li and Tsiakas~\cite{li2017} investigate differences between recessionary and expansionary periods, while Hammerschmid and Lohre~\cite{hammerschmid2018} construct macroeconomic regime factors to explain variation in aggregate stock return predictability. These distinctions are important because business-cycle regimes, macroeconomic regimes, and volatility regimes capture related but conceptually different dimensions of market conditions.

Although prior research establishes that predictive performance can vary across volatility or economic regimes, much of the regime-dependent equity premium literature focuses on a single short-term forecasting horizon. Conversely, multi-horizon studies generally examine how predictive performance changes with the horizon without explicitly conditioning the horizon--performance relationship on market volatility. This separation motivates a joint evaluation of forecast horizons and volatility regimes. In particular, it remains important to examine whether the pattern of predictive performance across horizons differs between low- and high-volatility market conditions.

\subsection{Temporal Stability and Hypothesis Development}

Equity return predictability may also evolve over time as economic conditions, market structures, monetary policies, and investor behavior change. Paye and Timmermann~\cite{paye2006} document substantial instability in return prediction models and identify structural changes in the relationship between predictors and future returns. Harvey et al.~\cite{harvey2021} further examine the real-time detection of regimes of predictability in the US equity premium, highlighting that periods of predictive ability may emerge and subsequently disappear. These findings suggest that aggregate OOS performance can conceal substantial variation across chronological subperiods.

Structural instability may also have implications for investors operating over different investment horizons. Pettenuzzo and Timmermann~\cite{pettenuzzo2011} show that accounting for structural breaks in return prediction models can materially alter long-horizon asset allocation decisions. Their findings emphasize that model instability is relevant not only to the estimation of predictive relationships but also to the interpretation of long-term investment outcomes. Nevertheless, structural-break analysis and chronological OOS evaluation address different aspects of temporal variation. The former identifies changes in model parameters or the underlying data-generating process, whereas the latter assesses whether predictive performance remains consistent across successive evaluation periods.

More importantly, temporal instability may affect not only overall predictive performance but also the relationship between forecast horizons and market volatility. For example, a forecasting model may perform relatively well at short horizons during low-volatility periods but lose its advantage during high-volatility periods. Such a pattern may also change across chronological subperiods. Therefore, examining temporal stability requires evaluating whether horizon-dependent and volatility-conditioned predictive patterns persist across different OOS evaluation windows, rather than relying solely on full-sample performance measures.

Building on the above literature, this study develops three hypotheses to examine the horizon dependence, volatility conditioning, and temporal stability of equity return predictability.

\begin{quote}
\noindent\textbf{H1:} \textit{Out-of-sample predictive performance differs significantly across forecast horizons.}

\medskip
\noindent\textbf{H2:} \textit{Out-of-sample predictive performance differs across volatility regimes, and the relationship between forecast horizons and predictive performance is conditional on market volatility.}

\medskip
\noindent\textbf{H3:} \textit{Out-of-sample predictive performance and its horizon--volatility regime patterns exhibit temporal instability across chronological subperiods.}
\end{quote}

Collectively, these hypotheses provide an empirical framework for investigating whether equity return predictability varies across investment horizons, whether such variation depends on prevailing market volatility, and whether the resulting patterns remain stable over time.


\section{Research Design and Methodology}\label{sec:data}

\subsection{Data Sources and Sample Construction}\label{subsec:data_sources}

This study uses daily financial data from the S\&P~500 index and the CBOE Volatility Index (VIX). Historical S\&P~500 data are obtained from Yahoo Finance~\cite{yahoosandp}, while VIX data are retrieved from the Federal Reserve Economic Data (FRED) database under the series identifier \texttt{VIXCLS}~\cite{fredvix}. The raw datasets contain 9,253 S\&P~500 observations (January 2, 1990--September 29, 2026) and 9,585 VIX observations (January 2, 1990--September 28, 2026), respectively.

The predictor variables are organized into six groups: lagged returns, momentum, historical volatility, VIX-related indicators, intraday indicators, and trading activity. These groups draw on prior research on return dynamics and volatility~\cite{french1987}, momentum~\cite{jegadeesh1993}, implied volatility and market uncertainty~\cite{banerjee2007,bekaert2014}, trading volume~\cite{campbell1993}, and range-based volatility estimation~\cite{alizadeh2002}. In total, the study employs 16 predictors, whose group-level composition is summarized in Table~\ref{tab:predictors}. Their exact definitions and construction procedures are provided in Section~\ref{subsec:targets_predictors}.

\begin{table}[htbp]
\caption{Summary of predictor variables across the six feature groups.}\label{tab:predictors}
\centering
\begin{tabular}{lcl}
\hline
Feature group & Number of predictors & Information captured \\
\hline
Lagged returns & 3 & Past return observations \\
Momentum & 3 & Cumulative return trends \\
Historical volatility & 3 & Recent return dispersion \\
VIX-related indicators & 4 & Implied market uncertainty \\
Intraday indicators & 2 & Intraday price movements \\
Trading activity & 1 & Changes in reported volume \\
\hline
Total & 16 & \\
\hline
\end{tabular}
\end{table}

Following Campbell, Lo, and MacKinlay~\cite{campbell1998}, continuously compounded returns (log returns) are defined as the natural logarithm of the ratio of consecutive asset prices. Since the S\&P~500 price index is used without incorporating dividend payments, the one-day log return is calculated as:
\begin{equation}\label{eq:daily_log_return}
R_t^{(1)} = \ln\left(\frac{P_t}{P_{t-1}}\right) = \ln(P_t) - \ln(P_{t-1}),
\end{equation}
where $P_t$ and $P_{t-1}$ denote the S\&P~500 closing index levels on trading days $t$ and $t-1$, respectively. The ``Close'' and ``Adj Close'' series are identical in the downloaded dataset. The first return observation is computed using the available open-to-close price movement but is subsequently excluded during feature construction.

The two datasets are aligned using the S\&P~500 trading calendar through a left join, with missing VIX observations forward-filled, resulting in 9,252 observations. Following the construction of forecasting targets and historical features, and the removal of observations with unavailable values, the final modeling dataset contains 9,073 observations from March~27, 1990, to April~7, 2026. The first 60 observations are excluded due to rolling-feature requirements, while the final 119 observations are unavailable for constructing the 120-day forward-return target.

\subsection{Forecasting Targets and Predictor Variables}\label{subsec:targets_predictors}

The study examines five cumulative forward-return horizons: 1, 5, 20, 60, and 120 trading days. For each horizon $h$, the forecasting target is defined as:
\begin{equation}\label{eq:target_return}
Y_t^{(h)} = \ln\left(\frac{P_{t+h}}{P_t}\right) = \sum_{i=1}^h R_{t+i}^{(1)}, \qquad h \in \{1, 5, 20, 60, 120\},
\end{equation}
where $P_t$ and $P_{t+h}$ denote the S\&P~500 closing index levels on trading days $t$ and $t+h$, respectively, and $R_{t+i}^{(1)}$ represents the one-day log return between trading days $t+i-1$ and $t+i$. Thus, the target measures the cumulative log return from the close of day $t$ to the close of day $t+h$.

\noindent\textbf{Predictor construction.} The 16 predictors are constructed from six groups of market information: lagged returns, momentum, historical volatility, VIX-related indicators, intraday price movements, and trading activity. Table~\ref{tab:predictor_definitions} presents the exact definitions and calculation procedures. All rolling statistics are calculated using observations available up to and including date $t$.

\begin{table}[htbp]
\centering
\caption{Definition and construction of predictor variables.}\label{tab:predictor_definitions}
\setlength{\tabcolsep}{6pt}
\renewcommand{\arraystretch}{1.2}
\begin{tabular}{lll}
\hline
Predictor group & Predictor & Definition \\
\hline
Lagged returns & \texttt{Return\_lag\_1} & $R_{t-1}^{(1)}$ \\
& \texttt{Return\_lag\_5} & $R_{t-5}^{(1)}$ \\
& \texttt{Return\_lag\_20} & $R_{t-20}^{(1)}$ \\
\hline
Momentum & \texttt{Momentum\_5D} & $\sum_{j=0}^{4} R_{t-j}^{(1)} = \ln(P_t/P_{t-5})$ \\
& \texttt{Momentum\_20D} & $\sum_{j=0}^{19} R_{t-j}^{(1)} = \ln(P_t/P_{t-20})$ \\
& \texttt{Momentum\_60D} & $\sum_{j=0}^{59} R_{t-j}^{(1)} = \ln(P_t/P_{t-60})$ \\
\hline
Historical volatility & \texttt{Volatility\_5D} & Sample standard deviation of $R_{t-4}^{(1)},\dots,R_t^{(1)}$ \\
& \texttt{Volatility\_20D} & Sample standard deviation of $R_{t-19}^{(1)},\dots,R_t^{(1)}$ \\
& \texttt{Volatility\_60D} & Sample standard deviation of $R_{t-59}^{(1)},\dots,R_t^{(1)}$ \\
\hline
VIX-related indicators & \texttt{VIXCLS} & $\text{VIX}_t$ \\
& \texttt{VIX\_lag\_1} & $\text{VIX}_{t-1}$ \\
& \texttt{VIX\_change\_1D} & $\text{VIX}_t - \text{VIX}_{t-1}$ \\
& \texttt{VIX\_pct\_change\_1D} & $(\text{VIX}_t - \text{VIX}_{t-1})/\text{VIX}_{t-1}$ \\
\hline
Intraday indicators & \texttt{HL\_range} & $(\text{High}_t - \text{Low}_t)/\text{Close}_t$ \\
& \texttt{OC\_return} & $(\text{Close}_t - \text{Open}_t)/\text{Open}_t$ \\
\hline
Trading activity & \texttt{Volume\_change} & $(\text{Volume}_t - \text{Volume}_{t-1})/\text{Volume}_{t-1}$ \\
\hline
\end{tabular}
\par\vspace{4pt}\parbox{\linewidth}{\fontsize{9pt}{12pt}\selectfont\textit{Note:} The rolling standard deviations used for the three historical-volatility predictors are calculated using the sample standard deviation (denominator $k-1$, consistent with the default pandas implementation). They are daily, non-annualized measures.}
\end{table}

The lagged-return predictors capture returns observed in previous trading sessions. Momentum predictors measure cumulative log returns over the corresponding rolling windows, while historical-volatility predictors measure the dispersion of daily log returns within those windows. In particular, the historical-volatility predictors are calculated using daily observations and are not annualized. Consistent with this convention, the 20-day rolling sample standard deviation (\texttt{Volatility\_20D}) is also used as the underlying non-annualized volatility measure for regime classification in Section~\ref{subsec:regime_classification}, where expanding historical quantiles convert it into ex-post evaluation regime labels that are excluded from the predictor matrix.

The four VIX-related predictors capture different aspects of option-implied market uncertainty. \texttt{VIXCLS} represents the VIX level associated with date $t$, while \texttt{VIX\_lag\_1} captures its previous-session level. \texttt{VIX\_change\_1D} measures the absolute daily change in the index, whereas \texttt{VIX\_pct\_change\_1D} measures its relative daily change.

The intraday indicators and trading-activity predictor are calculated as follows:
\begin{equation}\label{eq:hl_range}
\text{HL\_range}_t = \frac{\text{High}_t - \text{Low}_t}{\text{Close}_t},
\end{equation}
\begin{equation}\label{eq:oc_return}
\text{OC\_return}_t = \frac{\text{Close}_t - \text{Open}_t}{\text{Open}_t},
\end{equation}
\begin{equation}\label{eq:volume_change}
\text{Volume\_change}_t = \frac{\text{Volume}_t - \text{Volume}_{t-1}}{\text{Volume}_{t-1}}.
\end{equation}
Here, $\text{High}_t$, $\text{Low}_t$, $\text{Open}_t$, and $\text{Close}_t$ denote the corresponding daily S\&P~500 index levels, while $\text{Volume}_t$ denotes the volume value reported in the downloaded dataset. The high--low range captures the relative intraday price range, and the open--close return measures the simple price change during the trading session. Volume change is included as a proxy for market trading activity rather than as a direct measure of aggregate constituent-stock liquidity.

A zero-volume observation on May~24, 2023, is treated as a data anomaly. It is replaced with the preceding valid volume observation before calculating percentage changes. This procedure prevents an invalid zero denominator from affecting the volume-change feature and preserves the previous valid observation when calculating the subsequent change.

\noindent\textbf{Information availability and temporal alignment.} All 16 predictors are constructed using market observations associated with date $t$, without incorporating future observations. Four predictors---\texttt{Return\_lag\_1}, \texttt{Return\_lag\_5}, \texttt{Return\_lag\_20}, and \texttt{VIX\_lag\_1}---use information from previous trading sessions and are therefore available before the current trading session begins. The remaining 12 predictors incorporate observations associated with date $t$, including its completed price, volume, and VIX observations. Accordingly, the predictor vector is interpreted as a post-close information set rather than as information necessarily available before or during the closing auction.

Following the alignment procedure described in \hyperref[subsec:data_sources]{Section~\ref*{subsec:data_sources}}, the VIX series is matched to the S\&P~500 trading calendar using a left join, with missing observations forward-filled using the most recent preceding available value. This procedure preserves temporal consistency by carrying forward historical information without introducing future VIX observations. On dates requiring forward-filling, the VIX level remains unchanged, resulting in zero daily changes for the affected VIX-derived change measures.

The forecasting framework therefore evaluates the statistical predictability of subsequent cumulative returns conditional on information associated with the end of date $t$. Since some predictors, particularly current-day volume and the completed VIX observation, are only available after the relevant market observations have been finalized, the framework does not assume that a trading position can be executed at the observed closing price $P_t$. An executable trading strategy would require a subsequent execution opportunity and would consequently involve additional risks, including overnight price movements.

\subsection{Volatility Regime Classification}\label{subsec:regime_classification}

To examine regime-dependent predictability, observations are classified into Low, Normal, and High volatility regimes using the non-annualized 20-day rolling sample standard deviation of daily market log returns (\texttt{Volatility\_20D}):
\begin{equation}\label{eq:volatility}
\sigma_t^{(20)} = \operatorname{Std}_{\mathrm{sample}}\left(R_{t-19}^{(1)}, \dots, R_t^{(1)}\right).
\end{equation}
Expanding historical quantiles are employed to determine the volatility thresholds using information available up to time $t$~\cite{feng2024}:
\begin{equation}\label{eq:quantiles}
Q_{p,t} = \operatorname{Quantile}_p\left(\left\{\sigma_\tau^{(20)} : \tau \le t\right\}\right), \qquad p \in \{0.33, 0.67\}.
\end{equation}
A minimum of 252 historical observations is required before calculating the thresholds, with earlier observations back-filled using the first available expanding quantile thresholds. Observations are classified as Low when $\sigma_t^{(20)} \le Q_{0.33,t}$, Normal when $Q_{0.33,t} < \sigma_t^{(20)} \le Q_{0.67,t}$, and High when $\sigma_t^{(20)} > Q_{0.67,t}$.

Across the final sample, 3,062 observations (33.75\%) are classified as Low, 2,753 (30.34\%) as Normal, and 3,258 (35.91\%) as High volatility. Although the continuous underlying series $\sigma_t^{(20)}$ matches the 20-day historical-volatility predictor in Section~\ref{subsec:targets_predictors}, the discretized categorical regime labels are excluded from the predictor matrix and are not supplied to the forecasting models. All forecasting models are estimated globally using the same 16 predictors, with regime labels used exclusively for ex-post conditional evaluation.

The 20-day rolling window serves as the baseline specification for volatility-regime classification. To assess whether the regime-dependent findings are sensitive to the choice of rolling-window length, additional robustness analyses are conducted using 10-day and 60-day rolling standard deviations. These alternative specifications retain the same expanding-quantile classification procedure and are evaluated on the same date-aligned holdout sample. The corresponding results are reported in Section~\ref{subsec:volatility_regimes}.

\subsection{Chronological Data Splitting and Walk-Forward Validation}\label{subsec:data_splitting}

The dataset is divided chronologically into three non-overlapping subsets: training (6,351 observations; March~27, 1990--June~9, 2015), validation (1,360 observations; June~10, 2015--October~30, 2020), and hold-out testing (1,362 observations; November~2, 2020--April~7, 2026).

The validation period is evaluated through five consecutive windows of 272 observations using an expanding-window walk-forward procedure, following the rolling-origin framework of Tashman~\cite{tashman2000}. For each forecast horizon $h$, the final $h$ training observations are purged to ensure that training targets do not extend into the corresponding evaluation period:
\begin{equation}\label{eq:train_purged}
N_{\mathrm{train,purged}}^{(h)} = N_{\mathrm{train}} - h.
\end{equation}
At the transition to hold-out testing, the training sample is expanded to include both the original training and validation subsets, followed by horizon-specific purging. \texttt{StandardScaler} is fitted exclusively on the corresponding training data and then applied to validation or test observations, preventing information from the evaluation periods from influencing feature scaling, as emphasized by Hewamalage, Ackermann, and Bergmeir~\cite{hewamalage2023}. The same set of 16 predictors defined in Section~\ref{subsec:targets_predictors} is used consistently across the forecasting models and evaluation horizons.

For RQ3, the validation period is further divided into three non-overlapping calendar-based subperiods: 2015--2016 (395 observations), 2017--2018 (502 observations), and 2019--2020 (463 observations). The hold-out test period is separately organized into four subperiods: November--December 2020 (42 observations), 2021--2022 (503 observations), 2023--2024 (502 observations), and 2025--2026 (315 observations). Although the calendar labels overlap for 2019--2020, the observations are assigned according to the original chronological split boundaries and do not overlap between subsets.

The hold-out test set is reserved exclusively for final out-of-sample evaluation. Model configurations and hyperparameters are determined using training and validation data, with no model selection, hyperparameter optimization, or feature-engineering decisions based on hold-out performance. This design supports the evaluation of horizon-dependent predictability (RQ1), volatility-regime differences (RQ2), and temporal stability (RQ3) across five forecast horizons and three volatility regimes.


\section{Forecasting Models and Evaluation Framework}\label{sec:methodology}

\subsection{Forecasting Models and Benchmark}\label{subsec:models_benchmark}

This study evaluates twelve forecasting models across five forecast horizons: 1, 5, 20, 60, and 120 trading days. The models represent four methodological categories: econometric models (ARIMA and ARIMAX), penalized linear models (Ridge Regression and Elastic Net), kernel- and tree-based models (Support Vector Regression (SVR), Random Forest (RF), and XGBoost), and deep learning architectures (Long Short-Term Memory (LSTM), Temporal Convolutional Network (TCN), Vanilla Transformer, PatchTST, and TSMixer).

The selected models provide a comparative assessment of linear, nonlinear, and neural network-based approaches to equity return forecasting. ARIMA captures temporal dependence in the target return series, whereas ARIMAX incorporates additional explanatory information through exogenous predictors. Ridge Regression and Elastic Net examine the predictive contribution of regularized linear relationships. SVR, RF, and XGBoost provide kernel-based and tree-based alternatives for modeling nonlinear patterns. The deep learning models represent recurrent, convolutional, attention-based, patch-based, and feature-temporal mixing approaches to sequential forecasting.

The econometric models use the target return series and, in the case of ARIMAX, ten exogenous predictors: \textit{VIXCLS}, \textit{VIX\_lag\_1}, \textit{VIX\_change\_1D}, \textit{VIX\_pct\_change\_1D}, \textit{Momentum\_5D}, \textit{Momentum\_20D}, \textit{Volatility\_20D}, \textit{HL\_range}, \textit{OC\_return}, and \textit{Volume\_change}. The remaining ten machine learning and deep learning models use the sixteen predictors described in \hyperref[sec:data]{Section~\ref*{sec:data}}. A separate forecasting model is estimated for each forecast horizon.

The Historical Mean (HM) serves as the benchmark for evaluating out-of-sample forecasting performance~\cite{welch2008}. For each forecast horizon $h$, the benchmark forecast is calculated as the arithmetic mean of the eligible training targets:
\begin{equation}\label{eq:historical_mean}
\widehat{\mu}_{\mathrm{HM}}^{(h)}
=
\frac{1}{N_{\mathrm{train,purged}}^{(h)}}
\sum_{t\in \mathcal{T}_{\mathrm{train,purged}}^{(h)}}
Y_t^{(h)}.
\end{equation}
The HM benchmark is re-estimated independently for each validation fold and forecast horizon. During the final holdout evaluation, it is estimated using the combined purged training and validation sample. The estimated mean is positive across all evaluation windows, implying that HM consistently generates positive return forecasts within each window.

\subsection{Model Configuration and Hyperparameter Optimization}\label{subsec:model_config}

This subsection describes the model specifications, hyperparameter search spaces, optimization procedures, and training configurations used in the forecasting experiments.

\paragraph{Econometric model configurations.} ARIMA considers three candidate specifications: ARIMA(1,0,0), ARIMA(0,0,1), and ARIMA(1,0,1). The final specification is selected according to the Akaike Information Criterion (AIC), without differencing. ARIMAX adopts a fixed ARMA(1,1) specification and incorporates the ten exogenous predictors described in \hyperref[subsec:models_benchmark]{Section~\ref*{subsec:models_benchmark}}. The final specifications selected for each forecasting horizon are summarized in \hyperref[tab:config_econ_ml]{Table~\ref*{tab:config_econ_ml}}. For ARIMA, the selected order is determined independently for each horizon using the AIC criterion. ARIMAX retains the fixed ARMA(1,1) specification with the ten exogenous predictors across all horizons.

\paragraph{Conventional machine learning configurations.} Ridge Regression searches over the following regularization strengths:
\begin{equation}
\alpha \in \{10^{-3},10^{-2},10^{-1},1,10,100\}.
\end{equation}
Elastic Net jointly considers the regularization strength and the mixing parameter, with the search space defined as:
\begin{equation}
\begin{aligned}
\alpha &\in \operatorname{LogUniform}(10^{-5},1),\\
\ell_1\text{-ratio} &\in \{0.1,0.3,0.5,0.7,0.9\}.
\end{aligned}
\end{equation}
For SVR, the candidate hyperparameters include the penalty parameter ($C$), the kernel coefficient ($\gamma$), and the insensitive loss parameter ($\varepsilon$):
\begin{equation}
\begin{aligned}
C &\in \{0.1,1,10,100\},\\
\gamma &\in \{\text{scale},0.001,0.01,0.1\},
\qquad
\varepsilon \in \{0.001,0.01,0.05,0.1\}.
\end{aligned}
\end{equation}
For RF, the search space covers the number of estimators, maximum tree depth, minimum samples required to split an internal node, minimum samples required at a leaf node, and the number of features considered at each split. The candidate values include 200 and 500 estimators; maximum depths of unrestricted, 5, 10, and 20; minimum split sizes of 2, 5, and 10; minimum leaf sizes of 1, 2, and 5; and feature-selection strategies based on the square root of the available features or a proportion of 0.5. XGBoost is optimized over its corresponding tree-based hyperparameter space using the optimization procedure described below. The final hyperparameter configurations of the conventional machine learning models are reported in \hyperref[tab:config_econ_ml]{Tables~\ref*{tab:config_econ_ml}} and~\hyperref[tab:config_trees]{\ref*{tab:config_trees}}. These values represent the configurations selected for the final holdout experiments rather than the candidate search spaces used during optimization.

\begin{table}[htbp]
\centering
\caption{Final configurations of econometric and conventional regression models.}\label{tab:config_econ_ml}
\setlength{\tabcolsep}{4.5pt}
\begin{tabular}{llccccc}
\hline
Model & Hyperparameter & 1D & 5D & 20D & 60D & 120D \\
\hline
ARIMA       & Order               & (1,0,1) & (1,0,1) & (1,0,0) & (1,0,1) & (1,0,1) \\
ARIMAX      & Order               & (1,0,1) & (1,0,1) & (1,0,1) & (1,0,1) & (1,0,1) \\
ARIMAX      & Exogenous variables & 10      & 10      & 10      & 10      & 10 \\
Ridge       & Alpha               & 100     & 100     & 100     & 100     & 100 \\
Elastic Net & Alpha               & 0.001   & 0.001   & 0.005   & 0.05    & 0.05 \\
Elastic Net & L1 ratio            & 0.5     & 0.9     & 0.9     & 0.3     & 0.5 \\
SVR         & Kernel              & RBF     & RBF     & RBF     & RBF     & RBF \\
SVR         & $C$                 & 0.1     & 0.1     & 0.1     & 0.1     & 0.1 \\
SVR         & $\gamma$            & 0.001   & 0.001   & 0.001   & 0.001   & 0.1 \\
SVR         & $\varepsilon$       & 0.01    & 0.01    & 0.01    & 0.001   & 0.1 \\
\hline
\end{tabular}
\par\vspace{4pt}\parbox{\linewidth}{\fontsize{9pt}{12pt}\selectfont\textit{Notes:} ARIMA specifications are selected using AIC from three candidate models. ARIMAX uses ten exogenous predictors. Elastic Net uses \texttt{max\_iter} = 10,000. SVR uses the RBF kernel.}
\end{table}

\begin{table}[htbp]
\centering
\caption{Final configurations of tree-based machine learning models.}\label{tab:config_trees}
\setlength{\tabcolsep}{5pt}
\begin{tabular}{lrrrrr}
\hline
Hyperparameter & 1D & 5D & 20D & 60D & 120D \\
\hline
\multicolumn{6}{l}{\textbf{(a) Random Forest}} \\
\texttt{n\_estimators}      & 200    & 500    & 200    & 500    & 500 \\
\texttt{max\_depth}         & 5      & 5      & 5      & 5      & 5 \\
\texttt{min\_samples\_split}& 5      & 10     & 10     & 2      & 2 \\
\texttt{min\_samples\_leaf} & 1      & 2      & 2      & 2      & 5 \\
\texttt{max\_features}      & sqrt   & 0.5    & sqrt   & sqrt   & sqrt \\
\hline
\multicolumn{6}{l}{\textbf{(b) XGBoost}} \\
\texttt{n\_estimators}      & 200    & 500    & 300    & 900    & 500 \\
\texttt{learning\_rate}     & 0.0935 & 0.0449 & 0.1915 & 0.0902 & 0.0403 \\
\texttt{max\_depth}         & 9      & 5      & 9      & 3      & 3 \\
\texttt{min\_child\_weight} & 1      & 4      & 4      & 6      & 1 \\
\texttt{subsample}          & 0.722  & 0.694  & 0.642  & 0.663  & 0.772 \\
\texttt{colsample\_bytree}  & 0.656  & 0.780  & 0.625  & 0.807  & 0.932 \\
\texttt{gamma}              & 0.0015 & 0.0382 & 0.0990 & 4.2589 & 3.6227 \\
\texttt{reg\_alpha}         & 0.1411 & 0.0083 & 0.0002 & 0.0469 & 0.1260 \\
\texttt{reg\_lambda}        & 2.8797 & 0.1562 & 16.5550 & 0.2765 & 0.0171 \\
\hline
\end{tabular}
\par\vspace{4pt}\parbox{\linewidth}{\fontsize{9pt}{12pt}\selectfont\textit{Notes:} The reported values are the final configurations selected for the holdout evaluation. The hyperparameters are selected independently for each forecasting horizon.}
\end{table}

\paragraph{Hyperparameter optimization.} Hyperparameter optimization is conducted using two complementary strategies: grid-based search for conventional statistical and machine learning models, and Optuna-based optimization for selected tree-based and deep learning models.

For Optuna-based optimization, a Tree-structured Parzen Estimator (TPE) sampler with a fixed random seed of 42 is employed~\cite{akiba2019}. The optimization objective is to minimize the validation root mean squared error (RMSE) on an internal purged validation split, consisting of the final 15\% of the training data within each fold. This split is used exclusively for hyperparameter selection and is separated from the corresponding training subset to reduce temporal information leakage. The optimization budget is 30 trials for Random Forest (RF), 50 trials for XGBoost, and 30 trials each for LSTM and Vanilla Transformer.

The candidate search spaces for the four Optuna-optimized models are summarized in \hyperref[tab:optuna_search_spaces]{Table~\ref*{tab:optuna_search_spaces}}. For RF, the search space includes the number of estimators, maximum tree depth, minimum split and leaf sizes, and feature-selection strategy. XGBoost explores tree complexity, learning rate, sampling ratios, and regularization parameters. For LSTM and Vanilla Transformer, the search spaces additionally cover sequence length, architectural dimensions, dropout rates, learning rates, and batch sizes. During the Optuna stage, LSTM and Vanilla Transformer are trained for a maximum of 15 epochs per trial to limit computational costs. Their selected configurations are subsequently used in the final training procedure.

For the deep learning models, the \texttt{weight\_decay} parameter is not included in the actual Optuna search because it is not passed to the trial suggestion procedure in the corresponding model implementations. Accordingly, in the final holdout evaluation, LSTM and Vanilla Transformer use the Adam optimizer without explicitly specifying weight decay. The reported search spaces therefore reflect the hyperparameters actually explored by the implementation rather than all parameters declared in the configuration file.

The remaining models follow their respective predefined selection procedures. Ridge Regression, Elastic Net, and SVR use the candidate values specified in the conventional machine learning configurations above. ARIMA selects among three candidate specifications using the Akaike Information Criterion (AIC), while ARIMAX retains its predefined ARMA(1,1) specification with ten exogenous predictors. TCN, PatchTST, and TSMixer adopt predefined standardized benchmark configurations rather than undergoing the same Optuna-based optimization procedure. Consequently, these architectures do not receive an equivalent hyperparameter search budget to LSTM and Vanilla Transformer, a distinction that should be considered when interpreting comparisons among the deep learning models.

\begin{table}[htbp]
\centering
\caption{Hyperparameter search spaces and optimization budgets for Optuna-based models.}\label{tab:optuna_search_spaces}
\setlength{\tabcolsep}{4.5pt}
\begin{tabular}{lllr}
\hline
Model & Hyperparameter & Search space & Trials \\
\hline
RF & \texttt{n\_estimators}      & $\{200, 500\}$                    & 30 \\
   & \texttt{max\_depth}         & $\{\text{None}, 5, 10, 20\}$      & \\
   & \texttt{min\_samples\_split}& $\{2, 5, 10\}$                    & \\
   & \texttt{min\_samples\_leaf} & $\{1, 2, 5\}$                     & \\
   & \texttt{max\_features}      & $\{\text{sqrt}, 0.5\}$            & \\
\hline
XGBoost & \texttt{n\_estimators}     & $[100, 1000]$, step 100           & 50 \\
        & \texttt{learning\_rate}    & $\operatorname{LogUniform}(0.01, 0.2)$ & \\
        & \texttt{max\_depth}        & $[2, 10]$                         & \\
        & \texttt{min\_child\_weight}& $[1, 10]$                         & \\
        & \texttt{subsample}         & $[0.6, 1.0]$                      & \\
        & \texttt{colsample\_bytree} & $[0.6, 1.0]$                      & \\
        & \texttt{gamma}             & $[0, 5]$                          & \\
        & \texttt{reg\_alpha}        & $\operatorname{LogUniform}(10^{-8}, 10)$ & \\
        & \texttt{reg\_lambda}       & $\operatorname{LogUniform}(10^{-3}, 100)$ & \\
\hline
LSTM & \texttt{lookback}      & $\{10, 20, 40, 60\}$                        & 30 \\
     & \texttt{hidden\_size}   & $\{32, 64, 128\}$                           & \\
     & \texttt{num\_layers}    & $\{1, 2, 3\}$                               & \\
     & \texttt{dropout}       & $[0, 0.5]$                                  & \\
     & \texttt{learning\_rate} & $\operatorname{LogUniform}(10^{-5}, 10^{-3})$ & \\
     & \texttt{batch\_size}    & $\{16, 32, 64\}$                            & \\
\hline
Vanilla Transformer & \texttt{lookback}           & $\{20, 40, 60\}$                                   & 30 \\
                    & \texttt{d\_model}           & $\{32, 64\}$                                       & \\
                    & \texttt{nhead}              & $\{2, 4\}$                                         & \\
                    & \texttt{num\_encoder\_layers}& $\{1, 2\}$                                        & \\
                    & \texttt{dim\_feedforward}   & $\{64, 128\}$                                      & \\
                    & \texttt{dropout}            & $[0.05, 0.3]$                                      & \\
                    & \texttt{learning\_rate}      & $\operatorname{LogUniform}(10^{-5}, 5\times 10^{-4})$ & \\
                    & \texttt{batch\_size}         & $\{16, 32\}$                                       & \\
\hline
\end{tabular}
\par\vspace{4pt}\parbox{\linewidth}{\fontsize{9pt}{12pt}\selectfont\textit{Notes:} Optuna optimization uses a TPE sampler with seed 42 and minimizes validation RMSE on an internal purged split containing the final 15\% of each fold's training data. LSTM and Vanilla Transformer are trained for 15 epochs per trial. $\operatorname{LogUniform}$ denotes logarithmic sampling over the specified interval. Weight decay is excluded from the actual LSTM and Vanilla Transformer search spaces because it is not suggested by their respective Optuna implementations.}
\end{table}

\paragraph{Deep learning training configuration.} The deep learning architectures employ optimizer configurations according to their respective training implementations. LSTM and Vanilla Transformer are trained using the Adam optimizer without explicitly specifying weight decay, whereas TCN, PatchTST, and TSMixer use AdamW with a weight decay of $10^{-4}$~\cite{kingma2014}.

The learning rates and batch sizes are determined by the selected or predefined configurations for each model and forecasting horizon. LSTM uses horizon-specific batch sizes of 16 or 32, whereas Vanilla Transformer consistently uses a batch size of 16. TCN, PatchTST, and TSMixer use a learning rate of $10^{-4}$ and a batch size of 32 across all five forecasting horizons.

The maximum number of training epochs is set to 40 for LSTM and Vanilla Transformer and 1,000 for TCN, PatchTST, and TSMixer. Early stopping with a patience of seven epochs is applied based on validation MSE. Gradient clipping with a maximum norm of 1.0 is applied to improve optimization stability, and the models are trained using the MSE loss function. A fixed random seed of 42 is used to improve the reproducibility of the baseline training procedure. The complete final training configurations are summarized in \hyperref[tab:config_lstm_transformer]{Tables~\ref*{tab:config_lstm_transformer}} and~\hyperref[tab:config_dl_predefined]{\ref*{tab:config_dl_predefined}}.

\begin{table}[htbp]
\centering
\caption{Final configurations of LSTM and Vanilla Transformer.}\label{tab:config_lstm_transformer}
\setlength{\tabcolsep}{4.5pt}
\begin{tabular}{lrrrrr}
\hline
Hyperparameter & 1D & 5D & 20D & 60D & 120D \\
\hline
\multicolumn{6}{l}{\textbf{(a) LSTM}} \\
Lookback                & 20      & 40      & 40      & 10      & 40 \\
Hidden size             & 64      & 64      & 64      & 64      & 64 \\
Num.\ layers            & 3       & 1       & 3       & 3       & 3 \\
Dropout                 & 0.3539  & 0.4357  & 0.3861  & 0.4203  & 0.2136 \\
Learning rate           & 7.66e-4 & 4.05e-4 & 2.50e-5 & 5.34e-5 & 2.88e-4 \\
Batch size              & 16      & 32      & 32      & 32      & 16 \\
Optimizer               & Adam    & Adam    & Adam    & Adam    & Adam \\
Max.\ epochs            & 40      & 40      & 40      & 40      & 40 \\
Early stopping patience & 7       & 7       & 7       & 7       & 7 \\
\hline
\multicolumn{6}{l}{\textbf{(b) Vanilla Transformer}} \\
Lookback                & 60      & 40      & 20      & 20      & 60 \\
\texttt{d\_model}       & 32      & 32      & 64      & 64      & 64 \\
\texttt{nhead}          & 4       & 4       & 2       & 2       & 2 \\
\texttt{num\_layers}$^{\ast}$ & 2   & 2       & 2       & 2       & 2 \\
\texttt{dim\_feedforward} & 128   & 128     & 128     & 128     & 128 \\
Dropout                 & 0.2833  & 0.2963  & 0.2275  & 0.2767  & 0.1000 \\
Learning rate           & 2.02e-4 & 1.74e-4 & 4.80e-4 & 3.68e-4 & 4.86e-4 \\
Batch size              & 16      & 16      & 16      & 16      & 16 \\
Optimizer               & Adam    & Adam    & Adam    & Adam    & Adam \\
Max.\ epochs            & 40      & 40      & 40      & 40      & 40 \\
Early stopping patience & 7       & 7       & 7       & 7       & 7 \\
\hline
\end{tabular}
\par\vspace{4pt}\parbox{\linewidth}{\fontsize{9pt}{12pt}\selectfont\textit{Notes:} $^{\ast}$Important implementation note: although the tuning log stores the depth parameter as \texttt{num\_encoder\_layers} with values 1, 2, 2, 2, and 1, the holdout evaluation implementation retrieves \texttt{num\_layers} via \texttt{p.get('num\_layers', 2)}. Therefore, the effective number of encoder layers is 2 for all five horizons; the table reports the actual holdout configuration.}
\end{table}

\begin{table}[htbp]
\centering
\caption{Predefined configurations of TCN, PatchTST, and TSMixer.}\label{tab:config_dl_predefined}
\setlength{\tabcolsep}{6pt}
\begin{tabular}{lccc}
\hline
Hyperparameter & TCN & PatchTST & TSMixer \\
\hline
Lookback                & 20           & 20      & 20 \\
Architecture depth      & 3 levels     & 1 layer & 3 blocks \\
Hidden size / channels  & 64           & --      & -- \\
Number of channels      & [64, 64, 64] & --      & -- \\
Kernel size             & 3            & --      & -- \\
Patch length            & --           & 8       & -- \\
Stride                  & --           & 4       & -- \\
\texttt{d\_model}       & --           & 32      & -- \\
\texttt{nhead}          & --           & 2       & -- \\
\texttt{dim\_feedforward} & --         & 64      & -- \\
\texttt{ff\_dim}        & --           & --      & 64 \\
Dropout                 & 0.1          & 0.1     & 0.1 \\
Learning rate           & $10^{-4}$    & $10^{-4}$ & $10^{-4}$ \\
Batch size              & 32           & 32      & 32 \\
Optimizer               & AdamW        & AdamW   & AdamW \\
Weight decay            & $10^{-4}$    & $10^{-4}$ & $10^{-4}$ \\
Max.\ epochs            & 1,000        & 1,000   & 1,000 \\
Early stopping patience & 7            & 7       & 7 \\
\hline
\end{tabular}
\par\vspace{4pt}\parbox{\linewidth}{\fontsize{9pt}{12pt}\selectfont\textit{Notes:} These configurations are predefined and applied consistently across all five forecasting horizons. They are not selected through the same Optuna-based optimization procedure used for LSTM and Vanilla Transformer.}
\end{table}

\paragraph{Final configuration selection.} The final hyperparameter configurations for the holdout experiments are determined separately for each model and forecasting horizon using development-period validation results. For models undergoing fold-based hyperparameter optimization, the configuration from Fold 5 is prioritized because it corresponds to the longest training window. If a Fold 5 configuration is unavailable, the configuration with the lowest recorded \texttt{Best\_Val\_Score} is selected as a fallback. Models using fixed benchmark configurations or predetermined specifications retain those configurations without additional fold-based selection.

The selected configurations are frozen before the final holdout evaluation. No holdout observations are used for hyperparameter selection or model configuration adjustment. The complete configurations implemented in the holdout evaluation are reported in \hyperref[tab:config_econ_ml]{Tables~\ref*{tab:config_econ_ml}}--\hyperref[tab:config_dl_predefined]{\ref*{tab:config_dl_predefined}}.

\paragraph{Computational environment.} The computational experiments are conducted using Python 3.11.9 and PyTorch 2.6.0+cu124, with CUDA acceleration enabled on an NVIDIA GeForce RTX 4060 Laptop GPU with 8 GB of VRAM. The system has approximately 32 GB of RAM and runs Windows 11 Home. The principal software versions include scikit-learn 1.9.0, XGBoost 3.2.0, Optuna 5.0.0, statsmodels 0.15.0, and SciPy 1.17.1.

\paragraph{Multi-seed robustness evaluation.} To assess the sensitivity of deep learning forecasting performance to random initialization, an additional multi-seed evaluation is conducted on the final holdout set (Split 3: 2020--2026). The evaluation includes TSMixer with ten random seeds and LSTM, PatchTST, and Vanilla Transformer with five random seeds each. The seed values are specified before the experiments and are kept consistent across the evaluated architectures where applicable. The original model configurations and hyperparameters reported in \hyperref[tab:config_lstm_transformer]{Tables~\ref*{tab:config_lstm_transformer}} and~\hyperref[tab:config_dl_predefined]{\ref*{tab:config_dl_predefined}} are retained without further Optuna optimization or hyperparameter tuning. The training, validation, and holdout partitions, preprocessing procedures, forecasting horizons, and evaluation metrics remain unchanged across runs. For each model and horizon, a new model instance is initialized and trained independently under the corresponding random seed. This procedure isolates the sensitivity of forecasting performance to stochastic training initialization while controlling for changes in model configuration and data partitioning. The resulting out-of-sample $R^2$ values are summarized using the mean, standard deviation, median, minimum, maximum, and proportion of runs achieving positive OOS $R^2$. This supplementary evaluation is intended to assess initialization robustness rather than to select alternative model configurations or identify the best-performing random seed.

\subsection{Sequence Construction and Walk-Forward Implementation}\label{subsec:walk_forward}

The forecasting models are evaluated using the expanding-window walk-forward protocol and chronological data partitions described in \hyperref[subsec:data_splitting]{Section~\ref*{subsec:data_splitting}}. For sequence-based architectures, including LSTM, TCN, Vanilla Transformer, PatchTST, and TSMixer, historical observations are organized into input sequences according to the corresponding lookback length $L$.

When constructing validation and holdout inputs, the final $L-1$ observations from the preceding available data segment are retained as a historical lookback buffer. This allows the first observation in each evaluation segment to be forecast using the required historical context while preserving chronological ordering and avoiding the use of future information. The same sequence-construction procedure is applied consistently across validation folds and the final holdout evaluation.

\subsection{Forecast Evaluation Metrics}\label{subsec:eval_metrics}

Five metrics are used to evaluate forecasting performance: out-of-sample $R^2$, Root Mean Squared Error (RMSE), Mean Absolute Error (MAE), Directional Accuracy (DA), and Excess Directional Accuracy.

The primary metric is the out-of-sample $R^2$, which measures predictive performance relative to the Historical Mean (HM) benchmark \cite{campbell2008}:
\begin{equation}\label{eq:oos_r2}
R^2_{\mathrm{OOS}} = 1 - \frac{\sum_{t=1}^T (y_t^{(h)} - \hat{y}_t^{(h)})^2}{\sum_{t=1}^T (y_t^{(h)} - \hat{y}_{\mathrm{HM},t}^{(h)})^2}.
\end{equation}
Positive values indicate improvements over HM, whereas negative values indicate inferior performance. RMSE and MAE are used as supplementary measures of forecast error magnitude and are calculated according to their standard definitions \cite{hodson2022}.

Directional Accuracy (DA) measures the proportion of forecasts that correctly predict the sign of realized returns, following the directional forecasting evaluation framework discussed by Pesaran and Timmermann \cite{pesaran1992}:
\begin{equation}\label{eq:da}
\text{DA} = \frac{1}{T}\sum_{t=1}^T \mathbf{1}\left[\operatorname{sgn}(\hat{y}_t^{(h)}) = \operatorname{sgn}(y_t^{(h)})\right].
\end{equation}
To account for differences in the frequency of positive returns across forecast horizons, Excess Directional Accuracy is defined as:
\begin{equation}\label{eq:excess_da}
\text{Excess DA} = \text{DA} - \text{BaseRate}(h),
\end{equation}
where $\text{BaseRate}(h)$ denotes the proportion of positive realized returns at horizon $h$. Positive Excess DA indicates better directional accuracy than an always-positive prediction, whereas negative values indicate inferior performance.

\paragraph{Statistical inference for Excess DA.} To assess whether a candidate model achieves statistically significant directional improvement over the always-positive benchmark, Excess DA is evaluated using a one-sided Newey--West~\cite{newey1987} heteroskedasticity and autocorrelation consistent (HAC) test. For model $m$ and forecast horizon $h$, define the directional performance differential as:
\[
q_{m,t,h}
=
\mathbf{1}\left[
\operatorname{sgn}(\hat{y}_{m,t}^{(h)})
=
\operatorname{sgn}(y_t^{(h)})
\right]
-
\mathbf{1}\left[y_t^{(h)}>0\right].
\]
Its sample mean corresponds to Excess DA:
\[
\overline{q}_{m,h}
=
\frac{1}{T_h}\sum_{t=1}^{T_h}q_{m,t,h}
=
\mathrm{Excess\ DA}_{m,h},
\]
where $T_h$ denotes the number of valid out-of-sample observations at horizon $h$.

The null and one-sided alternative hypotheses are:
\[
H_0:E[q_{m,t,h}]\leq 0,
\qquad
H_1:E[q_{m,t,h}]>0.
\]
The test statistic is:
\[
Z_{m,h}
=
\frac{\overline{q}_{m,h}}
{\sqrt{\widehat{\Omega}_{m,h}/T_h}},
\]
where $\widehat{\Omega}_{m,h}$ is the Newey--West estimate of the long-run variance:
\[
\widehat{\Omega}_{m,h}
=
\widehat{\gamma}_0
+
2\sum_{j=1}^{L_h}
\left(1-\frac{j}{L_h+1}\right)
\widehat{\gamma}_j,
\]
with $\widehat{\gamma}_j$ denoting the sample autocovariance of the directional differential at lag $j$. The HAC truncation lag is matched to the forecast horizon, $L_h=h-1$, to account for the mechanical serial dependence induced by overlapping $h$-day cumulative return targets. The statistic is evaluated using a one-sided standard normal approximation.

The resulting $p$-values are adjusted using the Holm step-down procedure~\cite{holm1979} across all 60 model--horizon comparisons (12 candidate models and five forecast horizons). Statistical significance is assessed at the 5\% level using adjusted $p$-values.

Validation metrics are calculated as the arithmetic mean across the five walk-forward folds. Final performance is evaluated on the continuous hold-out sample of 1,362 observations. Results are further examined across forecast horizons, volatility regimes, and chronological subperiods.

As a benchmark robustness check, we additionally evaluate forecasting performance against a Zero-Return benchmark, which assumes zero expected cumulative returns over each forecast horizon, such that \(\hat{y}^{(h)}_{\mathrm{Zero},t}=0\). Unlike the primary Historical Mean benchmark, which uses the historical average of the target variable estimated from the training sample, the Zero-Return benchmark provides a fixed zero-return reference. Its corresponding out-of-sample \(R^2\) is calculated as:
\begin{equation}
R^2_{\mathrm{OOS,Zero}} = 1 - \frac{\sum_{t=1}^{T}\left(y_t^{(h)}-\hat{y}_t^{(h)}\right)^2}{\sum_{t=1}^{T}\left(y_t^{(h)}\right)^2}.
\end{equation}
The Historical Mean benchmark remains the primary reference for assessing predictive performance, while the Zero-Return benchmark is used as an alternative specification to examine the sensitivity of the results to benchmark selection. Both benchmarks are evaluated over the same hold-out observations for each model and forecast horizon. The resulting comparisons are reported in Section~\ref{subsec:benchmark_robustness}.

\subsection{Statistical Testing Framework}\label{subsec:hypothesis_testing}

The statistical testing framework examines differences in predictive performance across forecast horizons, volatility regimes, and evaluation periods. It comprises three main tests corresponding to H1--H3, supplemented by Diebold--Mariano (DM) tests and binomial sign tests. Newey--West heteroskedasticity and autocorrelation consistent (HAC) estimators are employed to account for serial dependence where applicable.

\paragraph{Normalized loss differentials.} Let $e_{\mathrm{HM},t,h}$ and $e_{m,t,h}$ denote the forecast errors of the historical mean (HM) benchmark and model $m$, respectively. The normalized squared error differential is defined as:
\begin{equation}
\widetilde d_{m,t,h}
=
\frac{e^2_{\mathrm{HM},t,h}-e^2_{m,t,h}}
{\mathrm{MSE}_{\mathrm{HM},h}},
\qquad
\mathrm{MSE}_{\mathrm{HM},h}
=
\frac{1}{T_h}\sum_{t=1}^{T_h}e^2_{\mathrm{HM},t,h}.
\end{equation}
A positive value indicates that model $m$ achieves a lower squared forecast error than the HM benchmark.

\paragraph{H1 -- Multi-horizon predictive performance.} H1 tests whether mean normalized predictive performance differs across the five forecast horizons $h\in\{1,5,20,60,120\}$. For each model, the horizon-specific differentials are arranged into a synchronized vector:
\[
X_{m,t}
=
\left(
\widetilde d_{m,t,1},
\widetilde d_{m,t,5},
\widetilde d_{m,t,20},
\widetilde d_{m,t,60},
\widetilde d_{m,t,120}
\right)^\top
\in\mathbb{R}^5,
\]
retaining only dates with valid observations at all five horizons. Let $\theta_m=E[X_{m,t}]$. The null hypothesis is:
\begin{equation}
H_0:R\theta_m=0,
\end{equation}
where
\begin{equation}
R=
\begin{bmatrix}
1&-1&0&0&0\\
0&1&-1&0&0\\
0&0&1&-1&0\\
0&0&0&1&-1
\end{bmatrix}.
\end{equation}
This restriction tests whether the five horizon-specific mean normalized loss differentials are equal. The Wald statistic is:
\begin{equation}
W_{\mathrm{H1},m}
=
T(R\bar X_m)^\top
\left(R\widehat\Omega_mR^\top\right)^{-1}
(R\bar X_m),
\end{equation}
where $\bar X_m$ is the sample mean vector, $T$ is the synchronized sample size, and $\widehat\Omega_m$ is the Newey--West long-run covariance matrix estimated jointly from the five-dimensional differential vector. The baseline HAC lag truncation is $L=20$, applied consistently across models. Under the null, the statistic asymptotically follows a chi-squared distribution. The degrees of freedom are determined by the effective rank of the restriction covariance matrix: four for eleven models and two for Elastic Net, whose differential series are degenerate at three horizons.

As a robustness check, the HAC bandwidth is varied over $L\in\{5,20,60,119\}$. In addition, a paired synchronized moving-block bootstrap is conducted using a fixed block length of $b=60$ observations and 500 replications. The same resampled time indices are applied across all five horizon-specific differential series, preserving their contemporaneous alignment. The bootstrap samples are recentered under the null, and studentized Wald statistics are recalculated using the baseline HAC bandwidth $L=20$. Bootstrap $p$-values are then adjusted for multiple testing using the Holm procedure.

\paragraph{H2 -- Volatility regimes and horizon--regime interaction.} H2 evaluates whether predictive ability depends on the prevailing volatility regime. Regimes are defined by the non-annualized 20-day rolling sample standard deviation of daily log returns, $\sigma_t^{(20)}$:
\[
\begin{aligned}
\mathrm{Low}:&\quad \sigma_t^{(20)}\leq Q_{0.33,t},\\
\mathrm{Normal}:&\quad Q_{0.33,t}<\sigma_t^{(20)}\leq Q_{0.67,t},\\
\mathrm{High}:&\quad \sigma_t^{(20)}>Q_{0.67,t},
\end{aligned}
\]
where quantiles are computed recursively using all observations up to time $t$ with an expanding window and a minimum of 252 observations to avoid look-ahead bias.

\textit{Single-horizon conditional predictive ability (CPA).} At each horizon $h$, the Giacomini--White~\cite{giacomini2006} conditional predictive ability (GW-CPA) test forms the three-dimensional moment vector:
\[
m_{m,t,h}=z_t\widetilde d_{m,t,h},
\]
using the regime instrument vector:
\begin{equation}
z_t=
\begin{bmatrix}
1\\
I(\mathrm{Low}_t)\\
I(\mathrm{High}_t)
\end{bmatrix},
\end{equation}
with Normal volatility serving as the reference regime absorbed into the intercept. The joint null hypothesis of zero conditional expected normalized loss differential across all three regimes is:
\begin{equation}
H_0:E[z_t\widetilde d_{m,t,h}]=0.
\end{equation}
Letting $\bar m_{m,h}=T^{-1}\sum_{t=1}^{T}z_t\widetilde d_{m,t,h}$, the Wald statistic is:
\begin{equation}
W_{\mathrm{CPA}}
=
\bar m_{m,h}^{\top}
\widehat V_{\mathrm{HAC}}^{-1}
\bar m_{m,h}
\overset{a}{\sim}\chi^2(3),
\end{equation}
where $\widehat V_{\mathrm{HAC}}$ is the Newey--West HAC covariance matrix of $\bar m_{m,h}$, estimated with a Bartlett kernel and fixed lag truncation $L=20$ across all horizons. For degenerate cases in which $\widetilde d_{m,t,h}\equiv0$ identically (Elastic Net at $h\in\{20,60,120\}$), the effective degrees of freedom reduce to zero ($\mathrm{df}=0$).

\textit{Horizon--regime interaction.} To test whether regime differentials vary across forecast horizons, observations are stacked across all five horizons $h\in H=\{1,5,20,60,120\}$ into the horizon-specific projection:
\begin{equation}
\widetilde d_{m,t,h}
=
\alpha_h+\beta_{h,L}I(\mathrm{Low}_t)
+\beta_{h,H}I(\mathrm{High}_t)+u_{m,t,h},
\end{equation}
where Normal volatility again serves as the reference regime. The null hypothesis tests whether the Low--Normal and High--Normal regime contrasts at horizons $h\in\{5,20,60,120\}$ differ from their corresponding contrasts at the 1-day baseline horizon:
\begin{equation}
H_0:
\begin{cases}
\beta_{h,L}-\beta_{1,L}=0,\\
\beta_{h,H}-\beta_{1,H}=0,
\end{cases}
\quad
\forall h\in\{5,20,60,120\}.
\end{equation}
The Wald statistic jointly tests these eight contrasts using the stacked five-horizon GMM/OLS covariance estimator:
\begin{equation}
W_{\mathrm{int}}
=
(R_{\mathrm{int}}\widehat\beta_{\mathrm{int}})^\top
\left(
R_{\mathrm{int}}\widehat V_{\mathrm{panel}}
R_{\mathrm{int}}^\top
\right)^{-1}
(R_{\mathrm{int}}\widehat\beta_{\mathrm{int}})
\overset{a}{\sim}\chi^2(8),
\end{equation}
where $\widehat V_{\mathrm{panel}}$ is estimated with Newey--West HAC standard errors using a Bartlett kernel and fixed lag truncation $L=20$, reducing to $\chi^2(4)$ for Elastic Net due to identical benchmark predictions at 20, 60, and 120 days.

\paragraph{H3 -- Temporal stability of predictive performance.} H3 evaluates whether predictive gains are stable over the evaluation sample or concentrated in particular calendar subperiods.

\textit{Giacomini--Rossi fluctuation test.} Following Giacomini and Rossi~\cite{giacomini2010}, for each model and horizon, a local standardized fluctuation statistic is computed over rolling windows of length $m_w=252$ trading days ($\tau=m_w,\ldots,T$, yielding 1,111 windows for $T=1,362$):
\begin{equation}
t_{m,\tau,h}
=
\frac{1}{\sqrt{m_w}}
\sum_{s=\tau-m_w+1}^{\tau}
\frac{\widetilde d_{m,s,h}}{\widehat\sigma_{\mathrm{HAC}}},
\qquad
\tau=m_w,\ldots,T,
\end{equation}
where $\widehat\sigma_{\mathrm{HAC}}$ is the Newey--West long-run standard deviation of $\widetilde d_{m,t,h}$, estimated over the full holdout sample using a Bartlett kernel with fixed lag truncation $L=20$ across all horizons. The summary statistic is the two-sided supremum:
\begin{equation}
S_{m,h}
=
\max_{\tau=m_w,\ldots,T}
|t_{m,\tau,h}|.
\end{equation}
To test $H_0:E[\widetilde d_{m,t,h}]=0$ for all $t$, empirical $p$-values are generated via a fixed-length moving-block bootstrap with 500 replications and block length $b=60$, applied to the demeaned differential series $\widetilde d_{m,t,h}-\bar{\widetilde d}_{m,h}$, holding $\widehat\sigma_{\mathrm{HAC}}$ fixed across bootstrap draws.

\textit{Subperiod stability test.} To complement the continuous fluctuation test, calendar block stability is evaluated using a block-level HAC Wald test. In the holdout sample, observations are partitioned into three non-overlapping calendar blocks: 2021--2022 ($N_1=503$), 2023--2024 ($N_2=502$), and 2025--2026 ($N_3=315$), covering 1,320 observations after excluding the 42 observations from November--December 2020. The same three-block framework is also applied to the validation sample across 2015--2016 (395), 2017--2018 (502), and 2019--2020 (463). Let $\bar{\widetilde d}_k$ denote the mean normalized squared error differential in block $k$, scaled by the full-sample HM MSE rather than within-block OOS $R^2$. The null hypothesis of equal block-mean predictive performance is:
\begin{equation}
H_0:
\bar{\widetilde d}_1
=
\bar{\widetilde d}_2
=
\bar{\widetilde d}_3.
\end{equation}
The Wald test statistic is:
\begin{equation}
W_{\mathrm{sub}}
=
(R\widehat\mu)^\top
(R\widehat\Sigma_{\mathrm{HAC}}R^\top)^{-1}
(R\widehat\mu)
\overset{a}{\sim}\chi^2(2),
\end{equation}
where $\widehat\mu=(\bar{\widetilde d}_1,\bar{\widetilde d}_2,\bar{\widetilde d}_3)^\top$, $R$ imposes the two linear equality restrictions, and $\widehat\Sigma_{\mathrm{HAC}}$ is estimated via a Newey--West Bartlett kernel with fixed lag truncation $L=20$ across all horizons. When a model's forecasts collapse identically to the HM benchmark across blocks (yielding zero differentials and $\mathrm{df}=0$, as observed for Elastic Net at longer horizons), the resulting statistic is identically zero.

\paragraph{Multiple-testing adjustment.} Across the hypothesis tests (H1, H2 CPA, H2 interaction, H3 fluctuation, and H3 subperiod), $p$-values are adjusted for multiple comparisons across the twelve models using the Holm step-down procedure. The adjustment is applied within each test family (e.g., across models within a given horizon for CPA, fluctuation, and subperiod block tests, and across models for joint tests). Both raw and adjusted $p$-values are reported; conclusions are drawn at the 5\% significance level based on adjusted values.

\paragraph{Supplementary predictive accuracy tests.} In addition to the hypothesis-specific tests, two conventional statistical tests are conducted on the full out-of-sample evaluation sample to provide complementary assessments of predictive accuracy.

The Diebold--Mariano (DM) test evaluates whether a candidate model achieves a significantly lower mean squared forecast error than the Historical Mean (HM) benchmark~\cite{diebold1995}. For model $m$ and forecast horizon $h$, the loss differential is defined as:
\[
d_{m,t,h}
=
e^2_{\mathrm{HM},t,h}-e^2_{m,t,h}.
\]
A positive value indicates a forecasting improvement over HM. The test statistic is calculated as:
\[
\mathrm{DM}_{m,h}
=
\frac{\bar d_{m,h}}
{\sqrt{\widehat{\operatorname{Var}}(\bar d_{m,h})}},
\]
with the long-run variance estimated using a rectangular spectral window truncated at lag $h-1$. The Harvey--Leybourne--Newbold (HLN) small-sample correction is then applied~\cite{buturac2022,diebold1995,harvey1997}:
\begin{equation}\label{eq:dm_stat}
\mathrm{DM}^{*}_{m,h}
=
\mathrm{DM}_{m,h}
\sqrt{
\frac{T+1-2h+h(h-1)/T}{T}
}.
\end{equation}
The corrected statistic is evaluated against a Student-$t$ distribution with $T-1$ degrees of freedom~\cite{buturac2022,harvey1997}. A one-sided alternative is adopted to test whether the candidate model significantly outperforms HM.

\paragraph{Supplementary directional accuracy test.} As a conventional reference, the exact one-sided Binomial Test is additionally reported to examine whether directional accuracy exceeds the 50\% random-guessing benchmark:
\[
H_0:p\leq 0.5,
\qquad
H_1:p>0.5.
\]
Let $K$ denote the number of correct directional predictions among $n$ non-zero realized returns. Under the independent Bernoulli assumption, the corresponding $p$-value is:
\[
p_{\mathrm{bin}}
=
P(K_{\mathrm{bin}}\geq K)
=
\sum_{j=K}^{n}
\binom{n}{j}(0.5)^n.
\]
However, the independence assumption underlying the conventional Binomial Test is generally violated when cumulative $h$-day return targets overlap. Consequently, successive directional outcomes may exhibit serial dependence, making the nominal Binomial $p$-values potentially unreliable. Therefore, this test is treated exclusively as a supplementary descriptive reference and is not used as the primary statistical evidence of directional forecasting skill.

The primary assessment of directional improvement is instead based on Excess DA and horizon-matched Newey--West HAC inference, as defined in \hyperref[subsec:eval_metrics]{Section~\ref*{subsec:eval_metrics}}. Unlike the Binomial Test, Excess DA evaluates directional improvement relative to the realized positive-return frequency, which corresponds to the directional accuracy of the always-positive HM benchmark.

Both the DM and supplementary Binomial tests are conducted on the full out-of-sample evaluation sample for all twelve candidate models across the five forecast horizons. No multiple-testing adjustment is applied to these supplementary tests; their $p$-values are interpreted as nominal individual significance levels. In contrast, Excess DA inference uses Holm-adjusted $p$-values across the 60 model--horizon comparisons.

\section{Empirical Results and Discussion}\label{sec:results}

\subsection{Out-of-Sample Forecast Performance}\label{subsec:oos_performance}

The out-of-sample evaluation covers 1,362 daily forecast origins from November 2, 2020, to April 7, 2026, across five forecast horizons (1D, 5D, 20D, 60D, and 120D). Twelve candidate models are compared against the Historical Mean (HM) benchmark, comprising seven econometric and conventional machine-learning models (ARIMA, ARIMAX, Ridge, Elastic Net, SVR, Random Forest, and XGBoost) and five deep-learning models (LSTM, Transformer, TCN, PatchTST, and TSMixer). All models are evaluated using the same observations and cumulative log-return targets. Forecast performance is assessed through OOS $R^2$, root mean squared error (RMSE), mean absolute error (MAE), and directional accuracy (DA).

\hyperref[tab:oos_r2_main]{Table~\ref*{tab:oos_r2_main}} reports OOS $R^2$ relative to HM. Conventional machine-learning models generally exhibit more stable performance at short and intermediate horizons, whereas deep-learning results are more heterogeneous. TSMixer records the strongest OOS $R^2$ at 20D (+3.80\%), 60D (+4.98\%), and 120D (+17.53\%). At the 120D horizon, only TSMixer and ARIMA achieve positive OOS $R^2$, at +17.53\% and +1.95\%, respectively. However, these point estimates should not be interpreted as evidence of statistically significant improvements. TSMixer's DM test is only nominally significant at 120D ($p=0.0485$) and does not remain significant after Holm correction, either within the 12-model horizon-specific family ($p_{\mathrm{Holm}}=0.5817$) or across all 60 model--horizon comparisons ($p_{\mathrm{Holm}}=1.0000$). Overall, no candidate model demonstrates statistically significant forecast improvement over HM after multiple-testing adjustment.

\begin{table}[htbp]
\centering
\caption{Out-of-sample $R^2$ by model and forecast horizon.}\label{tab:oos_r2_main}
\setlength{\tabcolsep}{4.5pt}
\begin{tabular}{lrrrrr}
\hline
Model & 1D & 5D & 20D & 60D & 120D \\
\hline
ARIMA & $-$0.49\% & $-$8.25\% & $-$12.70\% & $-$10.26\% & +1.95\% \\
ARIMAX & +1.18\% & $-$5.90\% & $-$4.61\% & +2.52\% & $-$0.35\% \\
Ridge & +1.50\% & +1.79\% & +1.66\% & +0.91\% & $-$4.95\% \\
Elastic Net & +0.13\% & +0.61\% & 0.00\% & 0.00\% & 0.00\% \\
SVR & +1.98\% & +1.48\% & +3.05\% & +4.64\% & $-$19.10\% \\
Random Forest & +1.77\% & $-$1.30\% & $-$0.39\% & +0.33\% & $-$6.42\% \\
XGBoost & +1.48\% & +0.11\% & +1.37\% & +0.13\% & $-$0.14\% \\
LSTM & $-$0.65\% & $-$0.50\% & +1.48\% & $-$0.96\% & $-$3.34\% \\
Transformer & $-$15.28\% & $-$39.72\% & $-$11.64\% & $-$21.28\% & $-$23.02\% \\
TCN & $-$14.16\% & $-$64.91\% & $-$123.43\% & $-$87.23\% & $-$142.55\% \\
PatchTST & $-$4.25\% & $-$6.85\% & $-$12.17\% & $-$21.86\% & $-$12.22\% \\
TSMixer & $-$2.72\% & $-$5.47\% & +3.80\% & +4.98\% & +17.53\% \\
\hline
\end{tabular}
\par\vspace{4pt}\parbox{0.85\linewidth}{\fontsize{9pt}{12pt}\selectfont\textit{Note:} OOS $R^2$ is calculated relative to the Historical Mean benchmark. Positive values indicate lower squared forecast error than HM, whereas negative values indicate underperformance.}
\end{table}

\FloatBarrier

\hyperref[tab:rmse_main]{Table~\ref*{tab:rmse_main}} and \hyperref[tab:mae_main]{Table~\ref*{tab:mae_main}} present RMSE and MAE, respectively. The results indicate that the model with the lowest error varies across forecast horizons and evaluation metrics. SVR achieves the lowest RMSE at 1D (0.0104), while TSMixer records the lowest RMSE at 20D (0.0409), 60D (0.0639), and 120D (0.0840). For MAE, SVR performs best at 5D (0.0164) and 60D (0.0489), whereas TSMixer achieves the lowest values at 20D (0.0308) and 120D (0.0663). At 120D, TSMixer improves upon HM in both RMSE (0.0840 versus 0.0925) and MAE (0.0663 versus 0.0757). In contrast, TCN exhibits substantially larger errors at longer horizons, with its 120D RMSE reaching 0.1441. These results highlight differences in point-forecast performance across model architectures, although lower error values alone do not establish statistically significant gains over the benchmark.

\begin{table}[htbp]
\centering
\caption{RMSE by model and forecast horizon.}\label{tab:rmse_main}
\setlength{\tabcolsep}{4.5pt}
\begin{tabular}{lrrrrr}
\hline
Model & 1D & 5D & 20D & 60D & 120D \\
\hline
HM & 0.0106 & 0.0223 & 0.0417 & 0.0655 & 0.0925 \\
ARIMA & 0.0106 & 0.0232 & 0.0442 & 0.0688 & 0.0916 \\
ARIMAX & 0.0105 & 0.0230 & 0.0426 & 0.0647 & 0.0927 \\
Ridge & 0.0105 & 0.0221 & 0.0413 & 0.0652 & 0.0948 \\
Elastic Net & 0.0105 & 0.0223 & 0.0417 & 0.0655 & 0.0925 \\
SVR & 0.0104 & 0.0222 & 0.0410 & 0.0640 & 0.1010 \\
Random Forest & 0.0105 & 0.0225 & 0.0418 & 0.0654 & 0.0954 \\
XGBoost & 0.0105 & 0.0223 & 0.0414 & 0.0655 & 0.0926 \\
LSTM & 0.0106 & 0.0224 & 0.0414 & 0.0658 & 0.0941 \\
Transformer & 0.0113 & 0.0264 & 0.0440 & 0.0722 & 0.1026 \\
TCN & 0.0113 & 0.0287 & 0.0623 & 0.0897 & 0.1441 \\
PatchTST & 0.0108 & 0.0231 & 0.0441 & 0.0723 & 0.0980 \\
TSMixer & 0.0107 & 0.0229 & 0.0409 & 0.0639 & 0.0840 \\
\hline
\end{tabular}
\par\vspace{4pt}\parbox{0.85\linewidth}{\fontsize{9pt}{12pt}\selectfont\textit{Note:} RMSE is measured in cumulative log-return units. Lower values indicate better point-forecast accuracy.}
\end{table}

\begin{table}[htbp]
\centering
\caption{MAE by model and forecast horizon.}\label{tab:mae_main}
\setlength{\tabcolsep}{4.5pt}
\begin{tabular}{lrrrrr}
\hline
Model & 1D & 5D & 20D & 60D & 120D \\
\hline
HM & 0.0075 & 0.0166 & 0.0327 & 0.0527 & 0.0757 \\
ARIMA & 0.0076 & 0.0171 & 0.0338 & 0.0532 & 0.0690 \\
ARIMAX & 0.0075 & 0.0172 & 0.0338 & 0.0512 & 0.0732 \\
Ridge & 0.0075 & 0.0165 & 0.0327 & 0.0514 & 0.0765 \\
Elastic Net & 0.0075 & 0.0165 & 0.0327 & 0.0527 & 0.0757 \\
SVR & 0.0075 & 0.0164 & 0.0316 & 0.0489 & 0.0867 \\
Random Forest & 0.0075 & 0.0167 & 0.0328 & 0.0534 & 0.0784 \\
XGBoost & 0.0075 & 0.0165 & 0.0327 & 0.0526 & 0.0759 \\
LSTM & 0.0076 & 0.0166 & 0.0322 & 0.0531 & 0.0760 \\
Transformer & 0.0084 & 0.0196 & 0.0364 & 0.0559 & 0.0849 \\
TCN & 0.0084 & 0.0209 & 0.0440 & 0.0649 & 0.0998 \\
PatchTST & 0.0078 & 0.0168 & 0.0344 & 0.0553 & 0.0759 \\
TSMixer & 0.0077 & 0.0169 & 0.0308 & 0.0497 & 0.0663 \\
\hline
\end{tabular}
\par\vspace{4pt}\parbox{0.85\linewidth}{\fontsize{9pt}{12pt}\selectfont\textit{Note:} MAE is measured in cumulative log-return units. Lower values indicate better point-forecast accuracy.}
\end{table}

Directional performance is evaluated using DA and Excess DA, defined as the difference between a candidate model's DA and that of HM, in percentage points (pp). \hyperref[tab:da_main]{Table~\ref*{tab:da_main}} reports both measures. Because HM issues positive forecasts throughout all horizons, its DA equals the corresponding positive-return base rate, increasing from 53.74\% at 1D to 78.05\% at 120D. Consequently, the increase in raw DA across longer horizons should not, by itself, be interpreted as an improvement in directional forecasting ability. Only TSMixer at 20D (+1.32 pp) and ARIMAX at 120D (+0.81 pp) numerically outperform HM. However, neither improvement is statistically significant under horizon-matched Newey--West HAC inference ($p=0.2758$ and $p=0.2035$, respectively). Furthermore, none of the 60 model--horizon comparisons demonstrates a statistically significant directional improvement after Holm correction. Nine comparisons produce exactly zero Excess DA because the corresponding models generate positive forecasts for all 1,362 observations, matching HM's directional predictions.

\begin{table}[htbp]
\centering
\caption{Directional accuracy and excess directional accuracy.}\label{tab:da_main}
{\small
\setlength{\tabcolsep}{5pt}
\begin{tabular}{lrrrrr@{\hspace{4pt}}c@{\hspace{4pt}}rrrrr}
\hline
 & \multicolumn{5}{c}{\textit{Panel A. DA (\%)}} & & \multicolumn{5}{c}{\textit{Panel B. Excess DA vs.\ HM (pp)}} \\
\cline{2-6}\cline{8-12}
Model & 1D & 5D & 20D & 60D & 120D & & 1D & 5D & 20D & 60D & 120D \\
\hline
HM & 53.74 & 59.40 & 66.89 & 74.96 & 78.05 & & -- & -- & -- & -- & -- \\
ARIMA & 50.37 & 53.01 & 59.10 & 68.21 & 76.87 & & $-$3.38 & $-$6.39 & $-$7.78 & $-$6.75 & $-$1.17 \\
ARIMAX & 52.94 & 50.00 & 53.45 & 72.10 & 78.85 & & $-$0.81 & $-$9.40 & $-$13.44 & $-$2.86 & \textbf{+0.81} \\
Ridge & 51.10 & 53.74 & 63.29 & 72.47 & 77.09 & & $-$2.64 & $-$5.65 & $-$3.60 & $-$2.50 & $-$0.95 \\
Elastic Net & 51.25 & 58.37 & 66.89 & 74.96 & 78.05 & & $-$2.50 & $-$1.03 & 0.00 & 0.00 & 0.00 \\
SVR & 52.72 & 58.08 & 66.81 & 74.96 & 65.12 & & $-$1.03 & $-$1.32 & $-$0.07 & 0.00 & $-$12.92 \\
Random Forest & 53.30 & 58.81 & 65.42 & 73.20 & 74.52 & & $-$0.44 & $-$0.59 & $-$1.47 & $-$1.76 & $-$3.52 \\
XGBoost & 53.16 & 59.25 & 66.52 & 74.96 & 78.05 & & $-$0.59 & $-$0.15 & $-$0.37 & 0.00 & 0.00 \\
LSTM & 50.07 & 51.69 & 66.89 & 74.96 & 78.05 & & $-$3.67 & $-$7.71 & 0.00 & 0.00 & 0.00 \\
Transformer & 48.53 & 51.98 & 42.14 & 62.56 & 72.03 & & $-$5.21 & $-$7.42 & $-$24.74 & $-$12.41 & $-$6.02 \\
TCN & 46.26 & 50.29 & 52.57 & 70.34 & 73.64 & & $-$7.49 & $-$9.10 & $-$14.32 & $-$4.63 & $-$4.41 \\
PatchTST & 50.88 & 56.98 & 50.88 & 74.30 & 74.89 & & $-$2.86 & $-$2.42 & $-$16.01 & $-$0.66 & $-$3.16 \\
TSMixer & 49.12 & 53.67 & 68.21 & 73.13 & 77.39 & & $-$4.63 & $-$5.73 & \textbf{+1.32} & $-$1.84 & $-$0.66 \\
\hline
\end{tabular}}
\par\vspace{4pt}\parbox{\linewidth}{\fontsize{9pt}{12pt}\selectfont\textit{Note:} Excess DA is calculated as candidate-model DA minus HM DA in percentage points (pp). Positive values indicate a numerical improvement over HM; they do not imply statistical significance. Values are rounded to two decimal places. The unrounded differences for TSMixer 20D and ARIMAX 120D are +1.3216 pp and +0.8076 pp, respectively.}
\end{table}

The comparison between magnitude-based and directional metrics reveals an important distinction. Although TSMixer achieves the highest OOS $R^2$ at 120D and lower RMSE and MAE than HM, its DA remains below the benchmark by 0.66 pp. Conversely, ARIMAX achieves a small positive directional difference at 120D (+0.81 pp), despite recording a slightly higher RMSE than HM and a negative OOS $R^2$ ($-$0.35\%). Thus, improvements in forecast magnitude accuracy do not necessarily translate into better directional predictions. Taken together, the results suggest that conventional machine-learning models provide relatively stable short- and intermediate-horizon forecasts, while TSMixer shows the strongest point-forecast performance at longer horizons. Nevertheless, the statistical evidence does not establish a robust improvement over HM in either forecast accuracy or directional performance after accounting for multiple comparisons.

\FloatBarrier
\subsection{Horizon-Dependent Predictability (RQ1)}

\noindent\textbf{Descriptive Evidence of Horizon-Dependent Predictability.} RQ1 examines how out-of-sample (OOS) predictive performance varies across forecast horizons. The aggregate hold-out results reveal a non-monotonic and model-specific pattern rather than a uniform decay in predictability.

At the shortest horizon (1D), SVR, Random Forest, and Ridge achieve positive OOS $R^2$ values of 1.98\%, 1.77\%, and 1.50\%, respectively. However, this pattern does not persist uniformly at 5D: Ridge and SVR retain positive OOS $R^2$ values of 1.79\% and 1.48\%, whereas Random Forest falls below the historical mean benchmark, with an OOS $R^2$ of $-$1.30\%. This distinction suggests that favorable short-horizon performance is model- and horizon-specific rather than consistent across the entire short-term range.

At intermediate horizons (20D--60D), SVR and TSMixer exhibit comparatively favorable results. SVR achieves OOS $R^2$ values of 3.05\% at 20D and 4.64\% at 60D, while TSMixer records 3.80\% and 4.98\%, respectively. Other models show more limited gains, including Ridge (1.66\% at 20D and 0.91\% at 60D), XGBoost (1.37\% and 0.13\%), and ARIMAX (2.52\% at 60D). These results indicate that predictive performance can remain positive at intermediate horizons, although the magnitude of the gains varies considerably across model classes.

A particularly notable result appears at the longest horizon (120D), where TSMixer achieves an OOS $R^2$ of 17.53\%. Its RMSE and MAE are 0.0840 and 0.0663, respectively, compared with 0.0925 and 0.0757 for the HM benchmark. This result suggests that longer-horizon predictability is not necessarily weaker and that certain architectures may capture patterns that are less evident at shorter horizons. Nevertheless, this finding should be interpreted as a descriptive result for the evaluated TSMixer implementation rather than evidence that longer horizons are generally more predictable.

The results also caution against assuming that greater model complexity necessarily leads to stronger predictive performance. Although some architectures achieve favorable results at particular horizons, the evaluated Transformer and TCN models do not consistently outperform simpler alternatives, while PatchTST produces negative OOS $R^2$ values across all five horizons. This pattern is consistent with the broader concerns in the literature that more complex forecasting architectures do not automatically translate into better financial prediction performance~\cite{welch2008},~\cite{gu2020}.

Importantly, the positive OOS $R^2$ values reported above are descriptive point estimates. They indicate that a model's aggregate squared forecast error is lower than that of the HM benchmark over the corresponding evaluation sample, but they do not by themselves establish statistically significant predictive superiority. As reported in \hyperref[subsec:oos_performance]{Section~\ref*{subsec:oos_performance}}, none of the model-versus-HM comparisons remains statistically significant after Holm correction. Therefore, the horizon-specific patterns should be interpreted as observed differences in predictive performance rather than conclusive evidence of reliable gains over the benchmark.

\medskip
\noindent\textbf{Statistical Evidence for H1.} To determine whether the observed differences across horizons are statistically distinguishable, H1 tests the joint null hypothesis that a model's mean normalized loss differentials are equal across all five forecast horizons. Because the normalization denominator is the horizon-specific HM mean squared error, the sample mean of each normalized differential is equivalent to the corresponding OOS $R^2$. Thus, H1 formally evaluates whether the model's relative predictive performance differs across horizons. It does not test whether the model significantly outperforms HM at any individual horizon, nor does it impose a monotonic-decay restriction.

Under the baseline HAC specification with bandwidth $L=20$, ARIMA is the only model for which the equality restriction is rejected after Holm correction ($W=15.7438$, $\mathrm{df}=4$, adjusted $p=0.0406$). This rejection reflects substantial variation in ARIMA's relative performance across horizons: its OOS $R^2$ is $-$0.49\% at 1D, $-$8.25\% at 5D, $-$12.70\% at 20D, $-$10.26\% at 60D, and 1.95\% at 120D. The result therefore indicates horizon-dependent performance, not strong predictive ability. In particular, rejecting equality across horizons should not be interpreted as evidence that ARIMA consistently outperforms the HM benchmark.

Other models show weaker statistical evidence. TSMixer, TCN, and SVR have unadjusted $p$-values below 0.05, but their Holm-adjusted $p$-values are 0.1620, 0.2520, and 0.4031, respectively. Consequently, their horizon-specific differences do not remain statistically significant after accounting for multiple comparisons.

The bandwidth sensitivity analysis provides additional evidence about the dependence of the results on the covariance specification. At $L=5$, five models---TCN, TSMixer, SVR, ARIMA, and Random Forest---remain significant after Holm correction. At $L=20$, $L=60$, and $L=119$, ARIMA is the only model that remains significant, with adjusted $p$-values of 0.0406, 0.0056, and 0.0108, respectively. Thus, ARIMA's rejection is consistent across the four examined bandwidths, whereas the evidence for the other models is sensitive to the selected HAC bandwidth.

The paired synchronized moving-block bootstrap provides a more conservative robustness check. Using a block length of $b=60$, 500 replications, and the baseline HAC bandwidth $L=20$, none of the twelve models remains significant after Holm correction. For ARIMA, the bootstrap $p$-value is 0.154, compared with the asymptotic HAC $p$-value of 0.0034. The bootstrap therefore does not corroborate the baseline asymptotic rejection. However, this result should also be interpreted cautiously: the block length of 60 observations is shorter than the maximum forecast horizon of 120 days, potentially limiting the preservation of longer-range serial dependence across resampled blocks, while 500 replications provide limited precision in estimating tail probabilities. Accordingly, the bootstrap results indicate inferential sensitivity rather than definitively invalidating the HAC findings.

Overall, the descriptive evidence reveals heterogeneous and non-monotonic horizon-dependent patterns, including favorable intermediate-horizon results for SVR and TSMixer and a pronounced 120D result for TSMixer. However, the formal H1 evidence is considerably more limited: only ARIMA rejects horizon-invariant performance under the baseline HAC specification, this result is not corroborated by the moving-block bootstrap, and the findings for other models are sensitive to the HAC bandwidth. These results do not establish a general or statistically robust monotonic decay in equity return predictability across forecast horizons.

\FloatBarrier
\subsection{Volatility Regime-Dependent Predictability (RQ2)}\label{subsec:volatility_regimes}

This section examines whether the out-of-sample (OOS) predictive performance of the evaluated models varies across volatility regimes and forecast horizons. The analysis distinguishes between descriptive regime-specific performance and formal statistical evidence, since differences in OOS $R^2$ across regimes do not necessarily imply statistically significant differences in predictive accuracy.

\medskip
\noindent\textbf{Regime-Specific Predictive Performance.}
Following the volatility classification procedure described in \hyperref[subsec:regime_classification]{Section~\ref*{subsec:regime_classification}}, the holdout evaluation sample contains 1,362 observations, comprising 332 Low-volatility, 599 Normal-volatility, and 431 High-volatility observations. These regime counts are identical across models and forecast horizons because all comparisons use the same date-aligned evaluation sample.

The historical-mean (HM) benchmark is estimated using the pre-holdout training and validation sample and remains fixed throughout the holdout period. Regime-specific OOS $R^2$ is evaluated relative to the HM benchmark within each corresponding volatility regime. Positive values indicate that a forecasting model outperforms HM within that regime, whereas negative values indicate inferior predictive performance.

\hyperref[tab:selected_regime_oos_r2]{Table~\ref*{tab:selected_regime_oos_r2}} provides a summary of descriptive predictive performance across volatility regimes and forecast horizons. It reports the distribution of positive, zero, and negative OOS $R^2$ outcomes, the number of models achieving positive OOS $R^2$ at each horizon, and selected model--horizon results illustrating cross-regime heterogeneity. To provide a more comprehensive view beyond these summary statistics, \hyperref[tab:full_regime_oos_r2]{Table~\ref*{tab:full_regime_oos_r2}} reports the complete regime-specific OOS $R^2$ results for all twelve forecasting models across all five forecast horizons. Together, the two tables provide both an aggregate overview and a detailed model-level comparison of predictive performance across volatility regimes.

Across the 60 model--horizon combinations within each regime, the Normal regime contains the largest number of positive OOS $R^2$ results, with 25 cases (41.7\%), compared with 16 cases (26.7\%) in the Low regime and 23 cases (38.3\%) in the High regime. As shown in the complete results in \hyperref[tab:full_regime_oos_r2]{Table~\ref*{tab:full_regime_oos_r2}}, negative OOS $R^2$ values nevertheless remain prevalent in all three regimes, indicating that most model--horizon combinations do not outperform the historical-mean benchmark. The detailed results also reveal substantial variation in the magnitude and direction of predictive performance across individual models and forecast horizons.

The distribution in \hyperref[tab:full_regime_oos_r2]{Table~\ref*{tab:full_regime_oos_r2}} also changes across forecast horizons. At the one-day horizon, seven models achieve positive OOS $R^2$ in the High regime, compared with six in the Normal regime and four in the Low regime. At the five-day horizon, only one model achieves a positive value in the Low regime, while five models do so in each of the Normal and High regimes. At the 60-day horizon, the pattern changes: six models are positive in the Low regime, five in the Normal regime, and only two in the High regime. This variation suggests that the relationship between volatility conditions and predictive performance is horizon-dependent rather than monotonic.

Several model-specific results in \hyperref[tab:full_regime_oos_r2]{Table~\ref*{tab:full_regime_oos_r2}} further illustrate this heterogeneity. TSMixer is a notable example of positive performance across all three regimes at the 120-day horizon, achieving OOS $R^2$ values of 17.81\%, 20.34\%, and 15.27\% in the Low, Normal, and High regimes, respectively. At the 20-day horizon, SVR also achieves positive OOS $R^2$ in all three regimes, with values of 0.12\%, 7.08\%, and 1.39\%, respectively. At the one-day horizon, Random Forest and XGBoost likewise produce positive values across all three regimes.

However, the complete results in \hyperref[tab:full_regime_oos_r2]{Table~\ref*{tab:full_regime_oos_r2}} show that positive performance at one horizon does not imply sustained predictive ability at longer horizons. SVR, for example, records positive OOS $R^2$ at the 60-day horizon in the Low and Normal regimes (3.49\% and 16.54\%, respectively), but its performance deteriorates at 120 days, reaching $-$33.02\% and $-$33.19\% in those regimes. Ridge exhibits a similar deterioration at 120 days, with OOS $R^2$ values of $-$17.17\%, $-$3.51\%, and $-$1.47\% across the Low, Normal, and High regimes.

As further illustrated by \hyperref[tab:full_regime_oos_r2]{Table~\ref*{tab:full_regime_oos_r2}}, the most pronounced adverse results occur for TCN under High volatility. Its OOS $R^2$ is $-$90.00\% at 5 days, $-$195.95\% at 20 days, $-$168.29\% at 60 days, and $-$289.23\% at 120 days. These results indicate substantial deterioration relative to the HM benchmark in this regime. To investigate the source of this deterioration, we further examine the distribution of TCN forecasts. At the 120-day horizon, the standard deviation of TCN predictions across the full holdout sample is 0.1053, compared with 0.0894 for realized returns, while the forecast--actual correlation is approximately zero ($-$0.005). The discrepancy becomes more pronounced under High volatility, where the forecast standard deviation reaches 0.1486, compared with 0.1119 for realized returns, and the mean predicted return increases to 0.1779. This combination of excessive forecast dispersion and negligible linear association with realized returns contributes to the substantial squared-error penalty relative to HM. For example, on April 21, 2022, TCN predicts a log return of 0.3937 when the realized value is $-$0.2056; on April 23, 2025, it predicts 0.7598 against a realized value of 0.2118. A similar pattern is observed at 20 days, where the forecast standard deviation is 1.16 times that of realized returns and High-regime OOS $R^2$ reaches $-$195.95\%.

This instability may be related to the model's architecture and training configuration. Unlike the regularized Ridge benchmark, which produces substantially less dispersed forecasts, the TCN implementation uses a fixed configuration with a dropout rate of 0.1 and residual convolutional blocks incorporating \texttt{BatchNorm1d}. Its final linear projection also imposes no explicit bound on forecast magnitude. These characteristics provide plausible explanations for the observed forecast dispersion, although the present results do not isolate the causal contribution of individual architectural components. Nevertheless, the evidence suggests that TCN's severe underperformance is associated not merely with reduced forecast accuracy under High volatility, but also with poorly calibrated forecast dispersion at longer horizons.

\begin{table}[htbp]
\centering
\caption{Summary of Regime-Specific Out-of-Sample Predictive Performance.}\label{tab:selected_regime_oos_r2}
\setlength{\tabcolsep}{6pt}
\begin{tabular}{llrrr}
\hline
\multicolumn{5}{l}{\textit{Panel A. Distribution of OOS $R^2$ outcomes}} \\
\hline
\multicolumn{2}{l}{Outcome} & Low & Normal & High \\
\hline
\multicolumn{2}{l}{Positive} & 16 (26.7\%) & 25 (41.7\%) & 23 (38.3\%) \\
\multicolumn{2}{l}{Zero} & 3 (5.0\%) & 3 (5.0\%) & 3 (5.0\%) \\
\multicolumn{2}{l}{Negative} & 41 (68.3\%) & 32 (53.3\%) & 34 (56.7\%) \\
\multicolumn{2}{l}{Total} & 60 (100\%) & 60 (100\%) & 60 (100\%) \\
\hline
\multicolumn{5}{l}{\textit{Panel B. Number of models with positive OOS $R^2$ by forecast horizon}} \\
\hline
\multicolumn{2}{l}{Forecast horizon} & Low & Normal & High \\
\hline
\multicolumn{2}{l}{1 day} & 4 & 6 & 7 \\
\multicolumn{2}{l}{5 days} & 1 & 5 & 5 \\
\multicolumn{2}{l}{20 days} & 1 & 4 & 6 \\
\multicolumn{2}{l}{60 days} & 6 & 5 & 2 \\
\multicolumn{2}{l}{120 days} & 4 & 5 & 3 \\
\hline
\multicolumn{5}{l}{\textit{Panel C. Selected model--horizon results across regimes}} \\
\hline
Model & Horizon & Low & Normal & High \\
\hline
Random Forest & 1 day & $+$0.07\% & $+$0.11\% & $+$2.85\% \\
XGBoost & 1 day & $+$0.05\% & $+$0.02\% & $+$2.42\% \\
SVR & 20 days & $+$0.12\% & $+$7.08\% & $+$1.39\% \\
TSMixer & 120 days & $+$17.81\% & $+$20.34\% & $+$15.27\% \\
SVR & 120 days & $-$33.02\% & $-$33.19\% & $-$2.98\% \\
TCN & 120 days & $-$26.22\% & $-$9.23\% & $-$289.23\% \\
\hline
\end{tabular}
\par\vspace{4pt}\parbox{\linewidth}{\fontsize{9pt}{12pt}\selectfont\textit{Notes:} OOS $R^2$ is expressed as a percentage. Panel~A classifies the 60 model--horizon combinations in each regime as positive, zero, or negative. Panel~B reports the number of models with positive OOS $R^2$ at each horizon. Panel~C presents selected cases to illustrate cross-regime heterogeneity and is not an exhaustive model comparison.}
\end{table}

For completeness, \hyperref[tab:full_regime_oos_r2]{Table~\ref*{tab:full_regime_oos_r2}} presents the full set of regime-specific OOS $R^2$ values for the twelve forecasting models across the five forecast horizons, reporting performance under the Low, Normal, and High volatility regimes for each model--horizon combination. This detailed breakdown allows direct comparison of individual model performance across volatility conditions and forecast horizons.

\begin{table}[htbp]
\centering
\caption{Full Regime-Specific Out-of-Sample $R^2$ (\%) Across All Models and Forecast Horizons.}\label{tab:full_regime_oos_r2}
{\footnotesize
\setlength{\tabcolsep}{4.5pt}
\begin{tabular}{lrrr@{\hspace{3pt}}c@{\hspace{3pt}}rrr@{\hspace{3pt}}c@{\hspace{3pt}}rrr}
\hline
\multicolumn{12}{l}{\textit{Panel A. Short- and intermediate-horizon regimes (1D, 5D, and 20D)}} \\
\hline
 & \multicolumn{3}{c}{1D} & & \multicolumn{3}{c}{5D} & & \multicolumn{3}{c}{20D} \\
\cline{2-4}\cline{6-8}\cline{10-12}
Model & Low & Norm. & High & & Low & Norm. & High & & Low & Norm. & High \\
\hline
ARIMA         & $+$0.05  & $-$0.47  & $-$0.59  & & $-$2.74  & $-$7.60  & $-$9.49  & & $-$6.67  & $-$5.99  & $-$17.15 \\
ARIMAX        & $-$0.01  & $+$0.85  & $+$1.55  & & $-$9.19  & $-$6.26  & $-$5.18  & & $-$4.38  & $-$8.76  & $-$2.47 \\
Ridge         & $-$1.06  & $+$0.53  & $+$2.41  & & $-$1.30  & $+$1.09  & $+$2.64  & & $-$1.95  & $+$0.52  & $+$2.81 \\
Elastic Net   & $+$0.15  & $+$0.11  & $+$0.13  & & $-$0.19  & $+$0.70  & $+$0.70  & & 0.00     & 0.00     & 0.00 \\
SVR           & $-$0.21  & $+$0.43  & $+$3.10  & & $-$1.18  & $+$0.89  & $+$2.21  & & $+$0.12  & $+$7.08  & $+$1.39 \\
Random Forest & $+$0.07  & $+$0.11  & $+$2.85  & & $-$0.07  & $+$0.22  & $-$2.24  & & $-$1.88  & $-$1.68  & $+$0.51 \\
XGBoost       & $+$0.05  & $+$0.02  & $+$2.42  & & $+$0.00  & $-$0.33  & $+$0.34  & & $-$0.38  & $-$0.99  & $+$2.88 \\
LSTM          & $-$2.70  & $-$1.72  & $+$0.22  & & $-$6.75  & $-$0.03  & $+$0.31  & & $-$0.12  & $+$3.88  & $+$0.46 \\
Transformer   & $-$22.78 & $-$29.70 & $-$7.20  & & $-$63.97 & $-$70.18 & $-$20.80 & & $-$38.53 & $-$29.49 & $+$1.86 \\
TCN           & $-$18.44 & $-$17.19 & $-$11.97 & & $-$39.39 & $-$22.23 & $-$90.00 & & $-$27.69 & $-$13.20 & $-$195.95 \\
PatchTST      & $-$0.74  & $-$4.60  & $-$4.71  & & $-$1.24  & $-$1.27  & $-$10.51 & & $-$14.62 & $-$8.48  & $-$13.73 \\
TSMixer       & $-$2.92  & $-$1.32  & $-$3.34  & & $-$2.74  & $+$2.20  & $-$9.67  & & $-$0.10  & $+$15.14 & $-$1.55 \\
\hline
\end{tabular}

\vspace{6pt}
\setlength{\tabcolsep}{5pt}
\begin{tabular}{lrrr@{\hspace{8pt}}c@{\hspace{8pt}}rrr}
\hline
\multicolumn{8}{l}{\textit{Panel B. Longer-horizon regimes (60D and 120D)}} \\
\hline
 & \multicolumn{3}{c}{60D} & & \multicolumn{3}{c}{120D} \\
\cline{2-4}\cline{6-8}
Model & Low & Norm. & High & & Low & Norm. & High \\
\hline
ARIMA         & $-$1.30  & $-$10.76 & $-$12.41  & & $-$7.43  & $+$11.61 & $-$1.99 \\
ARIMAX        & $-$0.53  & $-$4.13  & $+$8.09   & & $-$15.27 & $+$1.28  & $+$3.99 \\
Ridge         & $+$1.12  & $+$4.23  & $-$1.51   & & $-$17.17 & $-$3.51  & $-$1.47 \\
Elastic Net   & 0.00     & 0.00     & 0.00      & & 0.00     & 0.00     & 0.00 \\
SVR           & $+$3.49  & $+$16.54 & $-$3.49   & & $-$33.02 & $-$33.19 & $-$2.98 \\
Random Forest & $-$1.99  & $-$4.22  & $+$4.21   & & $+$1.10  & $-$11.74 & $-$5.12 \\
XGBoost       & $+$0.13  & $+$0.33  & $-$0.01   & & $-$0.41  & $-$0.83  & $+$0.49 \\
LSTM          & $+$0.71  & $-$2.32  & $-$0.46   & & $+$1.90  & $+$4.19  & $-$11.12 \\
Transformer   & $-$23.95 & $-$9.91  & $-$28.61  & & $-$21.17 & $-$28.80 & $-$19.26 \\
TCN           & $-$9.80  & $-$3.43  & $-$168.29 & & $-$26.22 & $-$9.23  & $-$289.23 \\
PatchTST      & $+$2.21  & $+$3.52  & $-$46.59  & & $+$8.29  & $+$0.57  & $-$29.79 \\
TSMixer       & $+$1.73  & $+$16.42 & $-$2.22   & & $+$17.81 & $+$20.34 & $+$15.27 \\
\hline
\end{tabular}}
\par\vspace{4pt}\parbox{\linewidth}{\fontsize{9pt}{12pt}\selectfont\textit{Notes:} Regime-specific OOS $R^2$ is expressed as a percentage relative to the Historical Mean (HM) benchmark within each volatility regime on the 1,362-observation holdout sample (Split~3: Low $n=332$, Normal $n=599$, High $n=431$). Positive values indicate lower squared forecast error than HM within the corresponding regime. For XGBoost at 5D under Low volatility, the unrounded OOS $R^2$ is strictly positive ($+0.0007\%$), rounding to $+0.00\%$ at two decimal places, whereas $0.00$ for Elastic Net at 20D, 60D, and 120D reflects forecasts identical to HM.}
\end{table}

\FloatBarrier
\noindent\textbf{Statistical Evaluation of Regime-Dependent Predictive Performance.}
Following the statistical framework established in \hyperref[subsec:hypothesis_testing]{Section~\ref*{subsec:hypothesis_testing}}, the Giacomini--White conditional predictive ability (GW-CPA) test is applied to evaluate whether forecasting performance differs significantly across the three volatility regimes. The tests use the normalized squared-error differentials and regime instruments defined in the methodological framework, with Newey--West HAC covariance estimation and Holm--Bonferroni correction applied across the twelve forecasting models separately within each horizon.

The results identify nine statistically significant GW-CPA rejections after Holm correction, as reported in \hyperref[tab:h2_gwcpa_sig]{Table~\ref*{tab:h2_gwcpa_sig}}. Importantly, all nine cases exhibit negative unconditional and regime-conditional mean differentials, indicating statistically significant underperformance relative to the HM benchmark rather than predictive gains.

The significant cases are concentrated in three models. TCN is significant at all five forecast horizons from 1 to 120 days. Transformer is significant at the 1-day and 5-day horizons, while ARIMA is significant at 5 days. SVR is significant at 120 days. No model demonstrates statistically significant outperformance over HM after the multiple-testing correction.

This distinction is particularly important for interpreting TSMixer at 120 days. Although TSMixer achieves positive OOS $R^2$ in all three regimes, its GW-CPA result does not remain statistically significant after Holm correction. Thus, the descriptive evidence of predictive gains should not be interpreted as statistically established superiority over the benchmark.

\begin{table}[htbp]
\centering
\caption{Statistically Significant GW-CPA Test Results After Holm Correction.}\label{tab:h2_gwcpa_sig}
\setlength{\tabcolsep}{4pt}
\begin{tabular}{llrrrrr}
\hline
Model & Horizon & Wald statistic & $\mathrm{df}$ & Raw $p$-value & Holm $p$-value & Mean differential \\
\hline
Transformer & 1D   & 30.9605 & 3 & 0.000001 & 0.000010 & $-$0.152796 \\
TCN         & 1D   & 40.6554 & 3 & $7.66\times 10^{-9}$ & $<$0.000001 & $-$0.141560 \\
ARIMA       & 5D   & 17.8814 & 3 & 0.000465 & 0.004653 & $-$0.082484 \\
Transformer & 5D   & 23.9464 & 3 & 0.000026 & 0.000308 & $-$0.397171 \\
TCN         & 5D   & 23.6924 & 3 & 0.000029 & 0.000319 & $-$0.649099 \\
TCN         & 20D  & 21.2389 & 3 & 0.000094 & 0.001127 & $-$1.234276 \\
TCN         & 60D  & 16.5454 & 3 & 0.000876 & 0.010517 & $-$0.872323 \\
SVR         & 120D & 13.0872 & 3 & 0.004452 & 0.048970 & $-$0.190970 \\
TCN         & 120D & 14.5199 & 3 & 0.002276 & 0.027318 & $-$1.425479 \\
\hline
\end{tabular}
\par\vspace{4pt}\parbox{\linewidth}{\fontsize{9pt}{12pt}\selectfont\textit{Notes:} The tests use three degrees of freedom for non-degenerate cases. Positive normalized squared-error differentials indicate improvement over HM, whereas negative values indicate underperformance. Holm--Bonferroni adjustment is conducted separately within each forecast horizon across the twelve models. All nine significant cases indicate underperformance relative to HM.}
\end{table}

\noindent\textbf{Horizon $\times$ Regime Interaction.}
While the GW-CPA test evaluates whether a model's predictive performance differs from the HM benchmark across volatility regimes, the horizon--regime interaction test examines whether the relative differences between regimes themselves change across forecast horizons.

The complete interaction-test results are presented in \hyperref[tab:h2_interaction]{Table~\ref*{tab:h2_interaction}}. TCN is the only model with a statistically significant horizon--regime interaction after Holm correction ($p_{\mathrm{Holm}}=0.002514$). Random Forest exhibits a relatively strong unadjusted result, but its adjusted $p$-value of 0.059990 does not meet the 5\% significance threshold. The remaining ten models do not show statistically significant interactions after correction.

The TCN results illustrate the economic magnitude of this interaction. Its normalized High--Normal differential gap changes from $-$0.1189 at the 1-day horizon to $-$1.5696 at 5 days and $-$4.1795 at 120 days. At the 120-day horizon, its regime-specific \texttt{Mean\_Skill} is $-$0.0755 in the Normal regime and $-$4.2550 in the High regime. These results indicate that TCN's relative performance deteriorates disproportionately under High volatility as the forecast horizon increases.

However, the interaction test should not be interpreted as a direct test of whether one regime is universally more predictable than another. It specifically tests whether the differences between regime-specific normalized loss differentials change across horizons.

\begin{table}[htbp]
\centering
\caption{Horizon $\times$ Regime Interaction Test Results.}\label{tab:h2_interaction}
\setlength{\tabcolsep}{6pt}
\begin{tabular}{lrrrr}
\hline
Model & Wald statistic & $\mathrm{df}$ & Raw $p$-value & Holm $p$-value \\
\hline
ARIMA         & 17.1180 & 8 & 0.028904 & 0.289037 \\
ARIMAX        &  9.2826 & 8 & 0.319019 & 1.000000 \\
Ridge         &  6.4439 & 8 & 0.597638 & 1.000000 \\
Elastic Net   &  3.4982 & 4 & 0.478153 & 1.000000 \\
SVR           &  9.4479 & 8 & 0.305929 & 1.000000 \\
Random Forest & 21.7240 & 8 & 0.005454 & 0.059990 \\
XGBoost       & 14.1973 & 8 & 0.076766 & 0.690897 \\
LSTM          & 12.1211 & 8 & 0.145880 & 1.000000 \\
Transformer   & 10.2816 & 8 & 0.245815 & 1.000000 \\
TCN           & 30.0223 & 8 & 0.000209 & 0.002514 \\
PatchTST      &  9.5971 & 8 & 0.294449 & 1.000000 \\
TSMixer       &  8.6626 & 8 & 0.371553 & 1.000000 \\
\hline
\end{tabular}
\par\vspace{4pt}\parbox{\linewidth}{\fontsize{9pt}{12pt}\selectfont\textit{Notes:} The interaction test examines whether the Low--Normal and High--Normal contrasts at the four non-baseline horizons differ from the corresponding contrasts at the 1-day horizon. The test has eight degrees of freedom for non-degenerate models. Elastic Net has four degrees of freedom because its forecast differentials are identically zero at the 20-, 60-, and 120-day horizons. Holm--Bonferroni correction is applied across the twelve models. Statistical significance is assessed at the 5\% level.}
\end{table}

\noindent\textbf{Distinguishing Regime-Specific OOS $R^2$ from \texttt{Mean\_Skill}.}
An important methodological distinction concerns the normalization used in regime-specific OOS $R^2$ and the \texttt{Mean\_Skill} statistic used in the GW-CPA analysis.

Regime-specific OOS $R^2$ uses the HM mean squared error calculated within the corresponding volatility regime. In contrast, \texttt{Mean\_Skill} uses the full-holdout HM mean squared error as a common normalization denominator for all regimes. Consequently, although both measures reflect the direction of relative squared-error performance, their magnitudes are not directly interchangeable because they use different benchmark normalization scales.

The relationship between the two measures has been verified across all 180 model--horizon--regime combinations. The difference arises because the regime-specific HM benchmark error can be substantially higher or lower than the full-holdout benchmark error.

This distinction is illustrated by TSMixer at the 120-day horizon. Its OOS $R^2$ is highest in the Normal regime (20.34\%), followed by Low (17.81\%) and High (15.27\%). However, its \texttt{Mean\_Skill} is highest in the High regime ($+$0.224579), followed by Normal ($+$0.166346) and Low ($+$0.127670). The difference arises because the High-regime HM MSE is approximately 1.471 times the full-holdout HM MSE, compared with ratios of approximately 0.818 in Normal and 0.717 in Low.

Accordingly, OOS $R^2$ is used to describe regime-specific predictive performance relative to the benchmark within each regime, whereas \texttt{Mean\_Skill} is used in the conditional predictive ability tests with a common full-holdout normalization. The two statistics serve different analytical purposes and should not be compared as though they were expressed on the same scale.

\medskip
\noindent\textbf{Overall Findings for RQ2.}
The findings provide evidence of substantial heterogeneity in predictive performance across volatility regimes and forecast horizons, but they do not support a universal relationship in which one volatility regime consistently favors forecasting models. The Normal regime contains the largest number of positive model--horizon results, yet individual models exhibit markedly different patterns. TSMixer demonstrates consistent positive performance across all three regimes at 120 days, whereas TCN experiences particularly severe deterioration under High volatility at longer horizons.

The formal statistical tests further qualify these descriptive findings. The GW-CPA results identify nine significant cases, all indicating underperformance relative to the historical-mean benchmark. Meanwhile, the horizon $\times$ regime interaction test identifies statistically significant variation in regime gaps only for TCN after multiple-testing correction.

Taken together, the results suggest that volatility conditions alone do not determine forecasting performance. Their association with predictive accuracy depends on the forecasting model and investment horizon. The evidence therefore supports a model- and horizon-dependent interpretation of volatility-regime predictability rather than a general claim that high or low volatility systematically improves equity return forecasts.

\subsubsection{Robustness to Alternative Volatility Window Definitions}\label{subsec:volatility_robustness}

To assess whether the findings for RQ2 depend on the 20-day rolling window used in the baseline regime classification, we repeat the regime-dependent evaluation using 10-day and 60-day rolling standard deviations. Both alternative specifications retain the expanding historical-quantile procedure described in Section~\ref{subsec:regime_classification} and use the same date-aligned holdout sample of 1,362 observations. The forecasting models, predictor sets, and evaluation framework remain unchanged.

The resulting regime distributions are broadly comparable across the three specifications, although the relative proportions of observations assigned to each regime vary with the window length. Under the 10-day specification, the holdout sample contains 354 Low-volatility observations (25.99\%), 567 Normal-volatility observations (41.63\%), and 441 High-volatility observations (32.38\%). The baseline 20-day specification contains 332, 599, and 431 observations in the respective regimes, while the 60-day specification contains 315, 662, and 385 observations. Thus, the longer window assigns a larger proportion of observations to the Normal regime and smaller proportions to both tail regimes.

Pairwise classification agreement further illustrates the sensitivity of regime assignments to the rolling-window length. The 10-day and 20-day specifications agree on 949 of 1,362 observations (69.68\%), while the agreement between the 20-day and 60-day specifications is 67.33\% (917 observations). Agreement between the two alternative specifications is lower, at 58.00\%. Importantly, direct switches between Low and High classifications are rare when either alternative is compared with the baseline: only 3 observations switch directly between these two regimes for the 10-day versus 20-day comparison, and 12 observations do so for the 20-day versus 60-day comparison. Most classification differences therefore involve reassignment between adjacent volatility regimes rather than a direct reversal between the two extremes.

Table~\ref{tab:volatility_window_robustness} summarizes the regime distributions and the main statistical results across the three rolling-window specifications. The number of statistically significant GW-CPA results after Holm correction is relatively similar, with 8, 9, and 6 significant model--horizon combinations under the 10-day, 20-day, and 60-day specifications, respectively. Under all three specifications, the significant cases indicate underperformance relative to the HM benchmark. The recurring results are concentrated in ARIMA at the 5-day horizon and Transformer at the 1-day and 5-day horizons. TCN also accounts for several significant cases, although the number of significant TCN horizons decreases under the 60-day specification. SVR at 120 days is significant under the 20-day and 60-day specifications, but not under the 10-day specification. Overall, the GW-CPA results do not indicate statistically significant predictive gains over HM under any of the three regime definitions.

\begin{table}[htbp]
\centering
\caption{Summary of Robustness Results Across Volatility Window Specifications.}\label{tab:volatility_window_robustness}
\setlength{\tabcolsep}{6pt}
\begin{tabular}{lccc}
\hline
\textbf{Measure} & \textbf{10-day} & \textbf{20-day} & \textbf{60-day} \\
\hline
\multicolumn{4}{l}{\textit{Regime distribution}} \\
Low volatility & 354 (25.99\%) & 332 (24.38\%) & 315 (23.13\%) \\
Normal volatility & 567 (41.63\%) & 599 (43.98\%) & 662 (48.60\%) \\
High volatility & 441 (32.38\%) & 431 (31.64\%) & 385 (28.27\%) \\
\hline
GW-CPA significant cases & 8/60 & 9/60 & 6/60 \\
Significant horizon--regime interactions & 0/12 & 1/12 & 2/12 \\
\hline
\end{tabular}
\par\vspace{4pt}\parbox{\linewidth}{\fontsize{9pt}{12pt}\selectfont\textit{Notes:} All specifications use the same 1,362-observation holdout sample. GW-CPA counts represent model--horizon combinations with statistically significant results after Holm correction. Interaction counts represent models with statistically significant horizon--regime interactions after Holm correction at the 5\% level. The 20-day specification is the baseline.}
\end{table}

The horizon--regime interaction tests provide a more differentiated picture. Under the 10-day specification, none of the twelve models exhibits a statistically significant interaction after Holm correction. Under the baseline 20-day specification, TCN is the only model with a significant interaction ($p_{\mathrm{Holm}}=0.0025$). Under the 60-day specification, both TCN ($p_{\mathrm{Holm}}=0.0324$) and Random Forest ($p_{\mathrm{Holm}}=0.0373$) exhibit significant interactions. Random Forest is close to, but does not reach, the 5\% significance threshold under the 10-day and 20-day specifications, with adjusted $p$-values of 0.0633 and 0.0600, respectively. TCN also approaches the threshold under the 10-day specification ($p_{\mathrm{Holm}}=0.0573$). These results suggest that evidence of horizon-dependent regime differences is sensitive to the volatility-window definition for some models, although TCN remains the most recurrent source of significant interaction effects.

The regime-specific OOS $R^2$ results provide additional descriptive evidence. Although the exact values and the number of positive model--horizon combinations vary across the three specifications, the main model-level patterns remain recognizable. TSMixer maintains positive OOS $R^2$ across all three regimes at the 120-day horizon under each window definition. Conversely, TCN continues to exhibit particularly severe deterioration under High volatility at longer horizons. These recurring patterns suggest that the principal descriptive findings are not solely artifacts of the baseline 20-day classification, even though the magnitude of regime-specific performance and the statistical evidence for horizon--regime interactions are not identical across specifications.

Taken together, the alternative-window analyses provide qualified support for the robustness of the RQ2 findings. The broad conclusion that predictive performance varies by model, forecast horizon, and volatility conditions remains consistent across the 10-day, 20-day, and 60-day specifications. However, the exact regime assignments and the statistical significance of individual horizon--regime interactions depend to some extent on the rolling-window length. Accordingly, the baseline 20-day specification remains the primary analysis, while the alternative windows demonstrate that the main conclusions are not entirely driven by this particular classification choice.

\FloatBarrier
\subsection{Temporal Stability of Predictive Performance (RQ3)}\label{subsec:temporal_stability}

RQ3 examines whether out-of-sample (OOS) predictive performance remains stable across calendar time or exhibits model- and horizon-specific differences. Temporal stability is evaluated using two complementary approaches: block-level HAC Wald tests and a Giacomini--Rossi-type fluctuation test with moving-block bootstrap inference.

The block-level analysis examines whether average predictive skill differs across predefined calendar periods, whereas the fluctuation test evaluates whether local predictive ability relative to the historical-mean (HM) benchmark differs from zero over rolling windows. These analyses focus on chronological variation and are conceptually distinct from the volatility-regime analysis in RQ2, which classifies observations into Low, Normal, and High volatility states.

\medskip
\noindent\textbf{Holdout Block Stability.}
Following the chronological sample partitioning described in \hyperref[subsec:data_splitting]{Section~\ref*{subsec:data_splitting}}, the holdout observations are divided into three non-overlapping calendar blocks: 2021--2022, 2023--2024, and 2025--2026. The 42 observations from November--December 2020 are excluded from this block-level analysis but retained in the aggregate holdout evaluation and rolling fluctuation test. Consequently, the block-level analysis covers 1,320 observations.

The HAC Wald tests examine whether mean normalized predictive skill is equal across the three calendar blocks. Predictive skill is measured using the normalized squared-error differential, with the full-holdout HM mean squared error serving as the common normalization denominator. Positive values indicate lower squared forecast errors than HM, whereas negative values indicate inferior predictive performance.

\hyperref[tab:h3_holdout_sig]{Table~\ref*{tab:h3_holdout_sig}} reports five model--horizon combinations with statistically significant differences across blocks after Holm correction: TCN at 20D, 60D, and 120D; LSTM at 60D; and PatchTST at 60D.

For TCN at 120D, mean predictive skill varies substantially across blocks ($-$3.6284, 0.0679, and $-$0.1418, respectively; adjusted $p=0.0047$), indicating considerable temporal heterogeneity.

In contrast, TSMixer at 120D shows no statistically significant block-level variation (Wald = 0.6653; adjusted $p=1.0000$). Its mean skill values are 0.2113, 0.0860, and 0.1472, with within-block OOS $R^2$ values of 15.90\%, 11.53\%, and 18.12\%, respectively. Although the mean skill values differ descriptively, the Wald test does not reject equality across blocks, suggesting comparatively consistent predictive performance over the three evaluation periods.

However, the absence of statistically significant block-level differences for TSMixer at 120D should be interpreted in conjunction with its sensitivity to random initialization. The block-level results reported above are based on the baseline configuration using seed 42 and therefore characterize temporal variation conditional on a single initialization. To assess whether the observed predictive performance is reproducible, TSMixer was additionally evaluated across ten independent random seeds with fixed hyperparameters. \hyperref[tab:tsmixer_multiseed]{Table~\ref*{tab:tsmixer_multiseed}} summarizes the resulting OOS $R^2$ distributions across all five forecast horizons.

The multi-seed results reveal substantial initialization sensitivity, particularly at longer forecast horizons. TSMixer achieves negative mean OOS $R^2$ at every horizon, ranging from $-$1.67\% at 1D to $-$21.49\% at 60D. Performance dispersion also increases with the forecast horizon, with the standard deviation rising from 0.98\% at 1D to 19.09\% at 120D. The proportion of runs achieving positive OOS $R^2$ remains low: none of the ten runs produces positive predictive skill at either 1D or 5D, while only two runs at 20D and one run at each of 60D and 120D outperform the historical-mean (HM) benchmark. These results indicate that the predictive performance observed under a particular initialization is not consistently reproduced across independent training runs.

The discrepancy is particularly evident at 120D. Although the seed-42 configuration achieves a positive aggregate OOS $R^2$ of 17.48\% and exhibits positive within-block OOS $R^2$ across all three holdout periods, the ten-seed evaluation produces a mean OOS $R^2$ of $-$18.93\% (SD = 19.09\%), with only one of ten runs achieving positive predictive skill. The median of $-$20.10\% further indicates that the negative mean is not driven solely by an isolated unfavorable run. Moreover, the mean OOS $R^2$ across the ten seeds is negative in all three calendar blocks ($-$6.18\%, $-$14.61\%, and $-$64.87\%, respectively). Thus, the favorable block-level performance observed under seed 42 does not represent the typical outcome across alternative initializations.

These findings qualify the interpretation of the block-level Wald test for TSMixer at 120D (Wald = 0.6653; adjusted $p = 1.0000$). The non-rejection indicates no statistically detectable difference in mean predictive skill across the three calendar blocks for the seed-42 configuration; it does not establish that TSMixer consistently maintains favorable predictive performance across repeated training runs. More broadly, \hyperref[tab:tsmixer_multiseed]{Table~\ref*{tab:tsmixer_multiseed}} demonstrates that temporal consistency conditional on a particular fitted model should be distinguished from reproducibility across stochastic training initializations. Consequently, the favorable temporal profile of TSMixer at 120D should be interpreted cautiously rather than as evidence of robust forecasting performance.

\begin{table}[htbp]
\centering
\caption{Multi-seed robustness of TSMixer across forecast horizons.}\label{tab:tsmixer_multiseed}
\setlength{\tabcolsep}{4.5pt}
\begin{tabular}{lrrrrrr}
\hline
Horizon & Mean OOS $R^2$ & SD & Median & Min & Max & Positive seeds \\
\hline
1D   & $-$1.67\%  & 0.98\%  & $-$1.55\%  & $-$3.52\%  & $-$0.28\% & 0/10 \\
5D   & $-$2.73\%  & 2.15\%  & $-$2.57\%  & $-$6.34\%  & $-$0.11\% & 0/10 \\
20D  & $-$6.96\%  & 10.10\% & $-$4.08\%  & $-$25.34\% & +4.53\%   & 2/10 \\
60D  & $-$21.49\% & 15.15\% & $-$23.45\% & $-$38.26\% & +4.91\%   & 1/10 \\
120D & $-$18.93\% & 19.09\% & $-$20.10\% & $-$44.09\% & +17.48\%  & 1/10 \\
\hline
\end{tabular}
\par\vspace{4pt}\parbox{\linewidth}{\fontsize{9pt}{12pt}\selectfont\textit{Notes:} OOS $R^2$ is expressed as a percentage. Positive seeds indicate the number of runs with OOS $R^2$ greater than zero. TSMixer was evaluated using ten random seeds with fixed hyperparameters.}
\end{table}

\begin{table}[htbp]
\centering
\caption{Significant Holdout Block HAC Wald Test Results.}\label{tab:h3_holdout_sig}
\setlength{\tabcolsep}{4.5pt}
\begin{tabular}{llrrrr}
\hline
Model & Horizon & Wald statistic & $\mathrm{df}$ & Raw $p$-value & Holm-adjusted $p$-value \\
\hline
TCN & 20D & 10.9722 & 2 & 0.0041 & 0.0497 \\
TCN & 60D & 17.0106 & 2 & 0.0002 & 0.0024 \\
TCN & 120D & 15.6699 & 2 & 0.0004 & 0.0047 \\
LSTM & 60D & 14.3866 & 2 & 0.0008 & 0.0083 \\
PatchTST & 60D & 10.6183 & 2 & 0.0049 & 0.0495 \\
\hline
\end{tabular}
\par\vspace{4pt}\parbox{\linewidth}{\fontsize{9pt}{12pt}\selectfont\textit{Notes:} Holm correction is applied separately within each forecasting horizon across the twelve competing models. Statistical significance is assessed at the 5\% level.}
\end{table}

\noindent\textbf{Validation Block Stability.}
The validation sample is divided into three calendar blocks: 2015--2016 (395 observations), 2017--2018 (502 observations), and 2019--2020 (463 observations). The HAC Wald tests examine equality of mean normalized predictive skill across these periods.

Although some model--horizon combinations exhibit nominal significance, none remains significant after Holm correction. The smallest adjusted $p$-value is 0.1184 for PatchTST at 120D. Thus, the validation results provide no statistically significant evidence of systematic differences in average predictive skill across the three calendar periods.

\medskip
\noindent\textbf{Rolling Predictive Performance: Giacomini--Rossi-Type Fluctuation Test.}
To complement the block-level analysis, the Giacomini--Rossi-type fluctuation test specified in \hyperref[subsec:hypothesis_testing]{Section~\ref*{subsec:hypothesis_testing}} is applied to examine temporal variations in local predictive ability relative to the HM benchmark.

The analysis uses rolling windows of 252 trading days and Newey--West HAC standardization with a fixed bandwidth of 20 observations. Statistical inference is conducted using 500 moving-block bootstrap replications with a block length of 60 observations. The supremum statistic is evaluated over all 1,111 rolling windows in the holdout sample, and Holm correction is applied across the twelve forecasting models separately for each forecast horizon.

\hyperref[tab:h3_fluctuation_sig]{Table~\ref*{tab:h3_fluctuation_sig}} reports the resulting fluctuation statistics and corresponding $p$-values after multiple-testing adjustment. The null hypothesis is rejected after Holm correction for Transformer at 1D, ARIMA at 20D, and TCN at 60D. These findings indicate that local predictive ability differs significantly from the HM benchmark in at least one 252-observation evaluation window.

Importantly, rejection does not necessarily imply that predictive performance alternates between superior and inferior outcomes over time. The rolling statistics provide further context for interpreting these results.

For Transformer at 1D, the rolling statistics remain negative across all 1,111 windows, ranging from $-$4.8926 to $-$0.2339, indicating persistently inferior predictive performance relative to HM. Its significant result therefore reflects sustained underperformance rather than alternating periods of relative predictive advantage.

The persistent underperformance of Transformer at 1D can be further understood from its forecast distribution. The standard deviation of its predictions is approximately 0.0042, only 0.40 times the standard deviation of realized returns (0.0105), while the forecast--actual correlation is just 0.031. Thus, despite producing relatively low-variance forecasts, Transformer captures almost none of the linear variation in realized returns. Its regime-specific OOS $R^2$ is negative across all three volatility regimes, reaching $-$22.78\% in Low, $-$29.70\% in Normal, and $-$7.20\% in High volatility. This indicates that the model's persistent rolling underperformance is not confined to a particular volatility state. Rather, its forecasts fail to provide sufficient information about realized return movements to improve upon the HM benchmark. The low forecast dispersion should therefore not be interpreted as evidence of effective regularization or reliable predictive behavior.

For ARIMA at 20D and TCN at 60D, the rolling statistics range from $-$3.3391 to 0.4237 and from $-$7.5441 to 0.1435, respectively. Negative statistics account for 88.39\% of ARIMA's windows and 77.77\% of TCN's windows, indicating that their local predictive performance is predominantly inferior to HM.

\begin{table}[htbp]
\centering
\caption{Significant Giacomini--Rossi-Type Fluctuation Test Results.}\label{tab:h3_fluctuation_sig}
\setlength{\tabcolsep}{5pt}
\begin{tabular}{llrrr}
\hline
Model & Horizon & Supremum statistic & Raw $p$-value & Holm-adjusted $p$-value \\
\hline
Transformer & 1D  & 4.8926 & $<$0.002 & $<$0.002 \\
ARIMA       & 20D & 3.3391 & 0.002    & 0.024 \\
TCN         & 60D & 7.5441 & 0.004    & 0.048 \\
\hline
\end{tabular}
\par\vspace{4pt}\parbox{\linewidth}{\fontsize{9pt}{12pt}\selectfont\textit{Notes:} Empirical $p$-values are based on 500 moving-block bootstrap replications with a block length of 60. The minimum non-zero $p$-value resolution is 0.002; a result of zero exceedances is reported as $p<0.002$. Holm correction is applied separately within each forecast horizon across the twelve models.}
\end{table}

Descriptively, TSMixer at 120D records an average 252-day rolling normalized predictive skill of 0.138 (SD = 0.174), compared with $-$1.670 (SD = 2.414) for TCN. These statistics are calculated as the mean and sample standard deviation of rolling normalized predictive skill across the 1,111 windows. They are not within-window OOS $R^2$ measures, as their normalization uses the full-holdout HM MSE rather than the HM MSE calculated separately within each rolling window. The substantially larger dispersion for TCN indicates pronounced temporal variation in its local predictive performance.

This rolling instability is consistent with the forecast-distribution diagnostics reported in \hyperref[subsec:volatility_regimes]{Section~\ref*{subsec:volatility_regimes}}. At the 120-day horizon, TCN exhibits a forecast standard deviation exceeding that of realized returns, while its forecast--actual correlation remains approximately zero. Consequently, its large rolling dispersion should be interpreted as fluctuations in forecast-error performance rather than evidence of alternating periods of strong predictive ability. In contrast, Transformer at 1D exhibits a different failure pattern: its forecast distribution is considerably less dispersed than realized returns, yet its predictive correlation is also negligible. The two architectures therefore demonstrate distinct forms of underperformance: excessive forecast dispersion and temporal instability for TCN, versus insufficient predictive information despite low forecast dispersion for Transformer.

The rolling predictive skill results for TSMixer at 120D should likewise be interpreted in light of the multi-seed findings. The reported average rolling skill of 0.138 (SD = 0.174) is obtained from the baseline seed-42 configuration and does not represent the average rolling performance across alternative random initializations. The supplementary multi-seed evaluation demonstrates that the favorable aggregate predictive performance of this configuration is not representative of the typical outcome across repeated runs. Specifically, the seed-42 result of +17.48\% OOS $R^2$ contrasts with the ten-seed mean of $-$18.93\%, while nine of the ten individual runs produce negative aggregate OOS $R^2$ values.

Consequently, the comparatively favorable rolling statistics reported for TSMixer should be regarded as conditional on the baseline initialization. They provide evidence about the temporal behavior of one trained configuration but do not establish that the architecture consistently maintains favorable local predictive skill across independent training runs. This distinction is important when interpreting rolling-window evidence alongside initialization sensitivity, as temporal stability within a single fitted model does not necessarily imply reproducibility across alternative model initializations.

\medskip
\noindent\textbf{Overall Assessment of RQ3.}
The results provide qualified support for H3, although the evidence depends on the temporal dimension being examined. Block-level HAC Wald tests identify significant differences in average predictive skill across calendar periods for five model--horizon combinations in the holdout sample, whereas validation block tests yield no significant results after Holm correction.

The Giacomini--Rossi-type fluctuation test further identifies three model--horizon combinations with statistically significant departures from equal local predictive ability relative to HM. However, these rejections should not be interpreted uniformly as alternating predictive performance or continuous instability. Transformer at 1D, for example, exhibits persistent underperformance across all rolling windows rather than alternating periods of superior and inferior performance.

Overall, temporal performance characteristics are model- and horizon-dependent rather than universal. The findings highlight the importance of evaluating forecasting models through both chronological block comparisons and local rolling-window assessments, rather than relying exclusively on aggregate OOS performance.

The supplementary multi-seed evaluation provides an additional qualification to these temporal findings. Although the baseline TSMixer configuration exhibits relatively consistent predictive skill across the three holdout calendar blocks and comparatively favorable rolling statistics at 120D, these characteristics are not robust to changes in random initialization. Across ten independent runs, TSMixer records a negative mean OOS $R^2$ of $-$18.93\% (SD = 19.09\%) at 120D, with positive predictive skill observed in only one run. The corresponding negative mean performance across all three calendar blocks further indicates that the favorable temporal profile of the baseline configuration is conditional on its initialization rather than a consistently reproducible property of the architecture.

Taken together, the evidence suggests that temporal variation and initialization sensitivity represent distinct dimensions of forecasting robustness. A model may exhibit relatively stable performance across chronological evaluation blocks under a particular initialization while remaining highly sensitive to stochastic training conditions. Therefore, the findings provide qualified support for H3: predictive performance exhibits model- and horizon-specific temporal variation, but the additional multi-seed evidence demonstrates that temporal consistency alone is insufficient to establish reliable out-of-sample forecasting performance. Evaluations of financial forecasting models should consequently consider both chronological stability and reproducibility across random initializations.


\subsection{Multi-Seed Robustness of Deep Learning Models}\label{subsec:multiseed_robustness}

To examine whether the forecasting performance of deep learning architectures is sensitive to stochastic training initialization, a supplementary multi-seed evaluation is conducted using the same holdout sample and model configurations as the baseline experiments. The evaluation retains the original hyperparameters without additional optimization and keeps the data partitions, preprocessing procedures, forecasting horizons, and evaluation metrics unchanged. Each run independently initializes and trains the corresponding model under a predefined random seed. This design allows performance variation across runs to be examined while controlling for changes in model configuration and data partitioning.

TSMixer is evaluated using ten random seeds, whereas LSTM, PatchTST, TCN, and Vanilla Transformer are each evaluated using five random seeds. The TSMixer results are reported separately in \hyperref[tab:tsmixer_multiseed]{Table~\ref*{tab:tsmixer_multiseed}} in Section~\ref{subsec:temporal_stability}. Table~\ref{tab:dl_multiseed_robustness} summarizes the corresponding multi-seed results for the other four deep learning architectures. Specifically, Panel~A reports the mean, standard deviation, and median OOS $R^2$, while Panel~B reports the extreme bounds (minimum and maximum OOS $R^2$) together with the mean MAE, RMSE, and directional accuracy (DA) across repeated runs.

\begin{table}[htbp]
\centering
\caption{Multi-seed robustness of deep learning models across forecast horizons.}\label{tab:dl_multiseed_robustness}
\renewcommand{\arraystretch}{0.96}
\setlength{\tabcolsep}{4.5pt}
\textbf{Panel A: Out-of-sample $R^2$ distribution across random seeds}\\[3pt]
\begin{tabular}{llrrr}
\hline
Model & Horizon & Mean OOS $R^2$ & SD & Median \\
\hline
LSTM & 1D   & $+$0.80\% & 1.85\% & $-$0.10\% \\
LSTM & 5D   & $+$0.19\% & 1.09\% & $+$0.55\% \\
LSTM & 20D  & $-$0.75\% & 2.03\% & $-$0.49\% \\
LSTM & 60D  & $-$2.14\% & 2.13\% & $-$1.64\% \\
LSTM & 120D & $-$5.36\% & 5.46\% & $-$3.13\% \\
\hline
PatchTST & 1D   & $-$8.98\% & 6.77\% & $-$6.05\% \\
PatchTST & 5D   & $-$10.03\% & 4.62\% & $-$8.81\% \\
PatchTST & 20D  & $-$23.97\% & 10.50\% & $-$26.36\% \\
PatchTST & 60D  & $-$23.78\% & 14.03\% & $-$27.21\% \\
PatchTST & 120D & $-$23.43\% & 12.84\% & $-$20.26\% \\
\hline
TCN & 1D   & $-$14.62\% & 4.77\% & $-$14.44\% \\
TCN & 5D   & $-$54.69\% & 19.34\% & $-$57.72\% \\
TCN & 20D  & $-$83.30\% & 54.60\% & $-$59.17\% \\
TCN & 60D  & $-$75.42\% & 51.96\% & $-$85.88\% \\
TCN & 120D & $-$83.54\% & 52.63\% & $-$70.73\% \\
\hline
Vanilla Transformer & 1D   & $-$18.32\% & 11.31\% & $-$19.61\% \\
Vanilla Transformer & 5D   & $-$26.26\% & 14.35\% & $-$21.10\% \\
Vanilla Transformer & 20D  & $-$7.83\% & 3.78\% & $-$8.93\% \\
Vanilla Transformer & 60D  & $-$25.60\% & 18.62\% & $-$18.90\% \\
Vanilla Transformer & 120D & $-$13.13\% & 18.44\% & $-$6.28\% \\
\hline
\end{tabular}

\vspace{8pt}
\setlength{\tabcolsep}{3.2pt}
\textbf{Panel B: Range bounds, forecast errors, and directional accuracy}\\[3pt]
\begin{tabular}{llrrrrr}
\hline
Model & Horizon & Min $R^2$ & Max $R^2$ & MAE & RMSE & DA \\
\hline
LSTM & 1D   & $-$0.86\% & $+$3.06\% & 0.0076 & 0.0105 & 49.16\% \\
LSTM & 5D   & $-$1.27\% & $+$1.37\% & 0.0165 & 0.0223 & 54.61\% \\
LSTM & 20D  & $-$3.95\% & $+$1.48\% & 0.0330 & 0.0418 & 64.45\% \\
LSTM & 60D  & $-$5.85\% & $-$0.51\% & 0.0536 & 0.0662 & 74.76\% \\
LSTM & 120D & $-$13.81\% & $+$0.47\% & 0.0759 & 0.0949 & 78.05\% \\
\hline
PatchTST & 1D   & $-$19.12\% & $-$2.24\% & 0.0080 & 0.0110 & 49.77\% \\
PatchTST & 5D   & $-$16.84\% & $-$5.74\% & 0.0172 & 0.0234 & 54.73\% \\
PatchTST & 20D  & $-$32.66\% & $-$7.19\% & 0.0352 & 0.0463 & 56.27\% \\
PatchTST & 60D  & $-$40.97\% & $-$7.42\% & 0.0548 & 0.0728 & 72.00\% \\
PatchTST & 120D & $-$41.36\% & $-$11.86\% & 0.0792 & 0.1027 & 74.80\% \\
\hline
TCN & 1D   & $-$21.83\% & $-$10.30\% & 0.0083 & 0.0113 & 48.21\% \\
TCN & 5D   & $-$79.12\% & $-$29.41\% & 0.0209 & 0.0277 & 47.75\% \\
TCN & 20D  & $-$162.60\% & $-$34.06\% & 0.0422 & 0.0559 & 48.68\% \\
TCN & 60D  & $-$129.96\% & $-$21.27\% & 0.0644 & 0.0860 & 64.95\% \\
TCN & 120D & $-$139.12\% & $-$26.02\% & 0.0906 & 0.1243 & 72.66\% \\
\hline
Vanilla Transformer & 1D   & $-$34.81\% & $-$6.24\% & 0.0085 & 0.0115 & 49.85\% \\
Vanilla Transformer & 5D   & $-$47.03\% & $-$11.52\% & 0.0189 & 0.0251 & 51.09\% \\
Vanilla Transformer & 20D  & $-$11.79\% & $-$3.75\% & 0.0341 & 0.0433 & 53.67\% \\
Vanilla Transformer & 60D  & $-$58.01\% & $-$12.12\% & 0.0563 & 0.0733 & 66.68\% \\
Vanilla Transformer & 120D & $-$43.34\% & $+$4.69\% & 0.0752 & 0.0982 & 74.90\% \\
\hline
\end{tabular}
\par\vspace{4pt}\parbox{\linewidth}{\fontsize{9pt}{12pt}\selectfont\textit{Notes:} OOS $R^2$ values in Panel A, minimum and maximum bounds in Panel B, and directional accuracy (DA) are expressed as percentages. SD denotes the standard deviation across random seeds. All metrics are computed across five independent random-seed training runs for each architecture and forecast horizon.}
\end{table}

The results reveal substantial differences in initialization robustness across architectures. LSTM exhibits the most favorable aggregate OOS $R^2$ performance among the four models, with slightly positive mean values at 1D and 5D (+0.80\% and +0.19\%, respectively). However, its mean OOS $R^2$ becomes negative at longer horizons, declining to $-$5.36\% at 120D. The standard deviation also increases from 1.85\% at 1D to 5.46\% at 120D, indicating greater variation across initializations at the longest horizon. The median and minimum values further show that the positive short-horizon means do not imply uniformly positive performance across individual runs.

PatchTST produces negative mean OOS $R^2$ at all five horizons. Its mean performance is particularly weak at 20D, 60D, and 120D, with values of $-$23.97\%, $-$23.78\%, and $-$23.43\%, respectively. The maximum OOS $R^2$ remains negative for every horizon, indicating that none of the five evaluated initializations achieves positive aggregate predictive skill relative to the HM benchmark. The standard deviations, especially at 60D and 120D, also indicate non-negligible variation in performance across runs. Thus, the underperformance of PatchTST is not attributable to a single unfavorable initialization in this experiment.

TCN exhibits the weakest and most variable results among the evaluated architectures. Its mean OOS $R^2$ is negative across all horizons, ranging from $-$14.62\% at 1D to $-$83.54\% at 120D. The deterioration is particularly pronounced at 5D, 20D, 60D, and 120D. At 20D, for example, the mean OOS $R^2$ is $-$83.30\%, while the standard deviation reaches 54.60\%. The minimum value of $-$162.60\% further illustrates the severity of the model's sensitivity to initialization at this horizon. Although the standard deviations are large, the maximum OOS $R^2$ remains negative at every horizon. This suggests that the observed weakness is systematic across the evaluated runs rather than being driven solely by one extreme seed.

Vanilla Transformer also records negative mean OOS $R^2$ at all horizons. Its weakest average performance occurs at 5D and 60D, at $-$26.26\% and $-$25.60\%, respectively. At 20D, the mean OOS $R^2$ improves to $-$7.83\%, accompanied by a comparatively smaller standard deviation of 3.78\%. Nevertheless, the maximum OOS $R^2$ remains negative at this horizon. At 120D, the maximum becomes positive (+4.69\%), but the mean remains negative ($-$13.13\%) and the standard deviation is relatively large (18.44\%). This indicates that a favorable result can occur under a particular initialization without representing the typical outcome across runs.

The error and directional metrics provide additional context for interpreting these results. Across the evaluated architectures, MAE and RMSE generally increase with the forecast horizon, reflecting the greater difficulty of predicting returns over longer periods. Directional accuracy also tends to increase at longer horizons, particularly for LSTM, PatchTST, and TCN. However, higher DA does not necessarily correspond to positive OOS $R^2$. For example, LSTM achieves a mean DA of 78.05\% at 120D while its mean OOS $R^2$ remains $-$5.36\%. This distinction indicates that directional performance and relative squared-error performance capture different aspects of forecasting quality and should not be treated as interchangeable measures.

The multi-seed results should also be interpreted alongside the TSMixer evidence in \hyperref[tab:tsmixer_multiseed]{Table~\ref*{tab:tsmixer_multiseed}}. TSMixer records negative mean OOS $R^2$ across all five horizons, despite the baseline seed-42 configuration producing a positive OOS $R^2$ of +17.48\% at 120D. Across ten runs, its mean at this horizon is $-$18.93\%, with only one run achieving positive OOS $R^2$. The contrast demonstrates that favorable results from a single initialization may not represent the typical performance of an architecture.

Overall, the multi-seed evaluation provides evidence that deep learning forecasting performance is sensitive to random initialization, with the magnitude and direction of this sensitivity varying across architectures and horizons. LSTM is comparatively more robust in terms of aggregate OOS $R^2$, whereas PatchTST, TCN, and Vanilla Transformer consistently produce negative mean OOS $R^2$ across the evaluated runs. These findings reinforce the importance of reporting repeated-run performance rather than relying exclusively on a single baseline initialization. They also qualify the interpretation of the temporal stability results in Section~\ref{subsec:temporal_stability}: stability within a particular fitted model does not necessarily imply reproducibility across independent training initializations.

\subsection{Feature Contribution Analysis: Ablation Study}\label{subsec:ablation_study}

To investigate the contribution of different predictor groups to forecasting performance, an ablation analysis is conducted using three models: XGBoost, Ridge, and LSTM. These models represent tree-based, linear, and deep learning approaches, respectively. The analysis considers six feature groups: Lagged Returns, Momentum, Historical Volatility, VIX, Intraday, and Trading Volume. The full specification contains 16 predictors, as described in Section~\ref{sec:data}.

For each model and forecasting horizon, the contribution of a feature group is assessed by comparing the OOS $R^2$ of the corresponding ablated specification with that of the full model. The difference is defined as:
\begin{equation}\label{eq:delta_oos_r2}
\Delta \mathrm{OOS}\,R^2_g = \mathrm{OOS}\,R^2_{\text{Ablated},g} - \mathrm{OOS}\,R^2_{\text{Full}},
\end{equation}
where $g$ denotes the feature group removed from the model. A negative value indicates that removing the group reduces OOS $R^2$, suggesting that the group contributes positively to predictive performance in the full specification. Conversely, a positive value indicates that the ablated model achieves higher OOS $R^2$ than the full model, suggesting that the removed group does not improve---and may reduce---performance under the evaluated configuration. Values close to zero indicate little measured difference between the two specifications.

\begin{table}[htbp]
\centering
\caption{Feature ablation results for XGBoost ($\Delta \mathrm{OOS}\,R^2$, percentage points).}\label{tab:ablation_xgboost}
{\footnotesize
\setlength{\tabcolsep}{6pt}
\begin{tabular}{lrrrrr}
\hline
Feature specification & 1D & 5D & 20D & 60D & 120D \\
\hline
Full Model & 0.00 & 0.00 & 0.00 & 0.00 & 0.00 \\
Without Lagged Returns & $-$0.32 & $+$0.06 & $+$0.67 & 0.00 & $-$0.05 \\
Without Momentum & $+$0.05 & $-$0.11 & $+$0.51 & 0.00 & $-$0.05 \\
Without Historical Volatility & $-$0.67 & $-$0.54 & $+$1.20 & 0.00 & $-$0.05 \\
Without VIX & $+$0.01 & $+$0.53 & $+$0.65 & 0.00 & $+$0.18 \\
Without Intraday & $-$1.20 & $+$0.09 & $+$0.75 & 0.00 & $-$0.05 \\
Without Trading Volume & $+$0.60 & $+$0.21 & $+$0.37 & 0.00 & 0.00 \\
\hline
\end{tabular}
}
\par\vspace{4pt}\parbox{\linewidth}{\fontsize{9pt}{12pt}\selectfont\textit{Notes:} Values are calculated as the OOS $R^2$ of the ablated specification minus that of the full model. Positive values indicate an improvement after removing the corresponding feature group; negative values indicate a reduction. All values are expressed in percentage points.}
\end{table}

The XGBoost results indicate that the contribution of feature groups varies across forecast horizons. At 1D, removing Intraday produces the largest negative change in OOS $R^2$ ($-$1.20 percentage points), followed by removing Historical Volatility ($-$0.67 percentage points) and Lagged Returns ($-$0.32 percentage points). These results suggest that intraday information, historical volatility, and lagged returns are useful for short-horizon prediction in this specification. In particular, the deterioration following the removal of Intraday is larger than that associated with any other feature group at 1D.

At 5D, removing Historical Volatility also reduces OOS $R^2$ by 0.54 percentage points, while removing Momentum produces a smaller negative change of 0.11 percentage points. In contrast, removing VIX or Trading Volume improves performance by 0.53 and 0.21 percentage points, respectively. At 20D, all six ablated specifications produce non-negative changes, ranging from $+$0.37 to $+$1.20 percentage points. The largest improvement occurs when Historical Volatility is removed ($+$1.20 percentage points), followed by Intraday ($+$0.75 percentage points) and Lagged Returns ($+$0.67 percentage points). This pattern suggests that the full set of predictors does not uniformly improve medium-horizon forecasting performance for XGBoost.

At 120D, the changes are comparatively small. Removing VIX improves OOS $R^2$ by 0.18 percentage points, whereas removing Lagged Returns, Momentum, or Historical Volatility produces a small deterioration of 0.05 percentage points. The removal of Trading Volume produces no reported change. For 60D, all six ablation results are reported as 0.00. This uniform pattern should be interpreted cautiously: it may reflect genuinely unchanged performance or rounding in the reported values. The underlying unrounded results should therefore be checked before concluding that these feature groups have no effect at this horizon.

\begin{table}[htbp]
\centering
\caption{Feature ablation results for Ridge ($\Delta \mathrm{OOS}\,R^2$, percentage points).}\label{tab:ablation_ridge}
{\footnotesize
\setlength{\tabcolsep}{6pt}
\begin{tabular}{lrrrrr}
\hline
Feature specification & 1D & 5D & 20D & 60D & 120D \\
\hline
Full Model & 0.00 & 0.00 & 0.00 & 0.00 & 0.00 \\
Without Lagged Returns & $+$0.10 & $-$0.04 & $-$0.03 & $-$0.01 & $+$0.05 \\
Without Momentum & $+$0.04 & $+$0.29 & $-$0.08 & $-$1.27 & $-$0.76 \\
Without Historical Volatility & $-$0.37 & $-$0.38 & $-$0.55 & $+$1.82 & $+$4.30 \\
Without VIX & $-$0.50 & $-$0.44 & $-$1.13 & $+$0.93 & $+$6.53 \\
Without Intraday & $-$0.07 & $-$0.20 & $-$0.07 & $-$0.25 & $-$0.02 \\
Without Trading Volume & $-$0.07 & 0.00 & 0.00 & 0.00 & $-$0.01 \\
\hline
\end{tabular}
}
\par\vspace{4pt}\parbox{\linewidth}{\fontsize{9pt}{12pt}\selectfont\textit{Notes:} Values are calculated as the OOS $R^2$ of the ablated specification minus that of the full model. Positive values indicate an improvement after removing the corresponding feature group; negative values indicate a reduction. All values are expressed in percentage points.}
\end{table}

Ridge exhibits a different feature-contribution pattern from XGBoost, particularly at longer horizons. At 1D, removing VIX reduces OOS $R^2$ by 0.50 percentage points, while removing Historical Volatility produces a reduction of 0.37 percentage points. These are the largest negative changes at this horizon, suggesting that both volatility-related feature groups contribute positively to the short-horizon performance of Ridge. Removing Intraday or Trading Volume also results in small negative changes, whereas removing Lagged Returns or Momentum produces slight improvements.

At 5D and 20D, the removal of VIX and Historical Volatility continues to reduce OOS $R^2$. The largest deterioration at 20D occurs when VIX is removed ($-$1.13 percentage points), followed by Historical Volatility ($-$0.55 percentage points). This pattern indicates that volatility-related information is particularly relevant to Ridge at the shorter and intermediate horizons considered in the experiment.

The results change substantially at longer horizons. At 60D, removing Historical Volatility improves OOS $R^2$ by 1.82 percentage points, while removing VIX improves it by 0.93 percentage points. At 120D, these improvements become considerably larger, reaching $+$4.30 percentage points without Historical Volatility and $+$6.53 percentage points without VIX. The latter is the largest positive change in the Ridge ablation results. Given that the full Ridge model has a negative OOS $R^2$ of $-$5.01\% at 120D, the positive changes indicate that the corresponding ablated specifications perform better than the full model, although the magnitude of the improvement should be interpreted relative to the weak full-model baseline.

Momentum also displays horizon-dependent effects. Removing it produces a small improvement at 1D and a larger improvement at 5D, but reduces OOS $R^2$ at 60D and 120D. By comparison, removing Intraday results in small negative changes across all five horizons, suggesting a relatively consistent positive contribution within the Ridge specification. The removal of Trading Volume produces changes close to zero throughout, indicating limited measured contribution to Ridge performance in this experiment.

\begin{table}[htbp]
\centering
\caption{Feature ablation results for LSTM ($\Delta \mathrm{OOS}\,R^2$, percentage points).}\label{tab:ablation_lstm}
{\footnotesize
\setlength{\tabcolsep}{6pt}
\begin{tabular}{lrrrrr}
\hline
Feature specification & 1D & 5D & 20D & 60D & 120D \\
\hline
Full Model & 0.00 & 0.00 & 0.00 & 0.00 & 0.00 \\
Without Lagged Returns & $+$1.28 & $+$0.92 & $-$2.95 & $+$3.09 & $+$3.94 \\
Without Momentum & $+$0.75 & $+$0.44 & $-$0.92 & $+$4.91 & $+$3.08 \\
Without Historical Volatility & $+$2.15 & $+$1.62 & $-$0.33 & $+$1.99 & $+$6.19 \\
Without VIX & $+$0.96 & $+$0.99 & $-$1.52 & $+$0.56 & $+$9.64 \\
Without Intraday & $+$1.57 & $+$2.01 & $-$0.60 & $+$2.08 & $+$6.44 \\
Without Trading Volume & $+$1.69 & $+$1.62 & $-$1.73 & $+$1.18 & $-$2.86 \\
\hline
\end{tabular}
}
\par\vspace{4pt}\parbox{\linewidth}{\fontsize{9pt}{12pt}\selectfont\textit{Notes:} Values are calculated as the OOS $R^2$ of the ablated specification minus that of the full model. Positive values indicate an improvement after removing the corresponding feature group; negative values indicate a reduction. All values are expressed in percentage points.}
\end{table}

\FloatBarrier
The LSTM ablation results show a pronounced horizon-dependent pattern. At 1D and 5D, removing any of the six feature groups improves OOS $R^2$. At 1D, the largest improvement is obtained by removing Historical Volatility ($+$2.15 percentage points), followed by Trading Volume ($+$1.69 percentage points) and Intraday ($+$1.57 percentage points). At 5D, removing Intraday yields the largest improvement ($+$2.01 percentage points), while removing Historical Volatility or Trading Volume improves OOS $R^2$ by 1.62 percentage points in each case. These results suggest that, for the evaluated LSTM configuration, the full predictor set does not provide the most favorable specification at the two shortest horizons.

The pattern differs at 20D. Removing any of the six feature groups reduces OOS $R^2$, with the largest deterioration occurring when Lagged Returns are removed ($-$2.95 percentage points). Removing Trading Volume and VIX also produces negative changes of $-$1.73 and $-$1.52 percentage points, respectively. This indicates that the feature groups contribute more positively to the LSTM's intermediate-horizon performance than they do at 1D and 5D.

At 60D, removing any feature group except Trading Volume improves OOS $R^2$. The largest improvement occurs when Momentum is removed ($+$4.91 percentage points), followed by Lagged Returns ($+$3.09 percentage points). Removing Historical Volatility, Intraday, and VIX also produces positive changes. At 120D, the improvements associated with removing VIX ($+$9.64 percentage points), Intraday ($+$6.44 percentage points), and Historical Volatility ($+$6.19 percentage points) are particularly pronounced. However, removing Trading Volume reduces OOS $R^2$ by 2.86 percentage points, indicating that this group retains a positive contribution at the longest horizon in the evaluated LSTM specification.

Taken together, the ablation results demonstrate that feature-group contributions are model- and horizon-dependent rather than uniform across forecasting configurations. XGBoost shows relatively small changes in most specifications, with more pronounced effects at 1D and 20D. Ridge exhibits substantial differences associated with VIX and Historical Volatility at longer horizons, particularly at 120D. LSTM displays the strongest variation across feature groups and horizons, with widespread improvements from feature removal at short and long horizons but deteriorations at 20D.

These findings do not imply that a feature group is universally useful or harmful. The measured contribution depends on the forecasting model, horizon, and other predictors retained in the specification. In addition, a positive $\Delta \mathrm{OOS}\,R^2$ should not automatically be interpreted as evidence that the removed variables are intrinsically uninformative; it indicates only that the corresponding ablated specification performs better under the evaluated experimental setup. Similarly, negative changes indicate that the removed group contributes to the performance of the full specification, but do not establish a causal relationship.

Overall, the ablation analysis provides complementary evidence to the horizon- and regime-dependent results in Sections~\ref{subsec:oos_performance} and~\ref{subsec:volatility_regimes}. It suggests that the effectiveness of predictor groups depends not only on the forecast horizon but also on the model architecture. Consequently, the results support evaluating predictor contributions at the model--horizon level rather than assuming that a single feature configuration is optimal across all forecasting tasks.

\FloatBarrier

\subsection{Robustness to Alternative Benchmark Specifications: Zero-Return Benchmark}\label{subsec:benchmark_robustness}

To examine whether the evaluation of predictive performance is sensitive to the choice of benchmark, we conduct an additional robustness analysis using a Zero-Return benchmark alongside the primary Historical Mean (HM) benchmark. The Zero-Return benchmark assumes that the expected cumulative return is zero at each forecast horizon, whereas the HM benchmark uses the historical average of the target variable estimated from the training sample. Both benchmarks are evaluated using the same hold-out observations for each model--horizon pair.

\paragraph{Comparison of out-of-sample performance across benchmarks.} Table~\ref{tab:benchmark_comparison} reports the out-of-sample \(R^2\) values under the two benchmark specifications. Panel A presents results relative to HM, while Panel B reports results relative to Zero-Return.

Under the primary HM benchmark, predictive performance varies considerably across models and forecast horizons. Several conventional machine learning models achieve positive \(R^2_{\mathrm{OOS}}\) values at shorter and intermediate horizons. Ridge and SVR, for example, exhibit positive values across most horizons up to 60D. However, performance deteriorates for several models at longer horizons. Transformer and TCN produce substantially negative \(R^2_{\mathrm{OOS}}\) values across the evaluation horizons, indicating that their forecast errors exceed those of the HM benchmark. In contrast, TSMixer achieves positive values at 20D, 60D, and 120D, with its strongest result at 120D (17.535\%).

The results change considerably when the Zero-Return benchmark is adopted. More model--horizon combinations exhibit positive \(R^2_{\mathrm{OOS}}\), particularly at 60D and 120D. At 120D, for example, TSMixer, ARIMA, and Elastic Net achieve \(R^2_{\mathrm{OOS}}\) values of 38.631\%, 27.033\%, and 25.582\%, respectively. Nevertheless, these values measure performance relative to a zero-return forecast and should not be interpreted as evidence that these models outperform the more demanding HM benchmark. The comparison illustrates that conclusions about predictive performance can depend materially on the benchmark specification.

\paragraph{Differences between benchmark-based evaluations.} Table~\ref{tab:benchmark_difference} reports the difference between the two \(R^2_{\mathrm{OOS}}\) measures. The difference is positive for all 60 model--horizon combinations. Averaged across models, the gap increases from 0.205 percentage points at 1D to 1.173 at 5D, 5.291 at 20D, 16.733 at 60D, and 29.687 at 120D. The overall average difference is 10.618 percentage points.

This widening gap at longer horizons indicates that the choice of benchmark has a particularly substantial effect when evaluating cumulative returns. Importantly, the difference does not represent an improvement in the forecasts themselves. It reflects the change in the benchmark error used in the denominator of the out-of-sample \(R^2\) calculation. Consequently, a positive \(R^2_{\mathrm{OOS}}\) against Zero-Return does not necessarily imply positive predictive performance against HM. The HM benchmark is therefore retained as the primary reference throughout the study.

\paragraph{Statistical comparison.} The Diebold--Mariano tests with the Harvey--Leybourne--Newbold correction provide additional evidence on whether individual models significantly outperform each benchmark (summarized in Table~\ref{tab:dm_benchmark_summary}). At the reported 5\% significance level, 1 of the 60 model--horizon comparisons is statistically significant against HM (1.7\%), corresponding to TSMixer at 120D (\(p=0.0485\)). Against Zero-Return, 8 of the 60 comparisons are statistically significant (13.3\%). These include ARIMAX at 60D, Ridge at 5D and 20D, Random Forest at 60D, XGBoost at 20D, LSTM at 60D, TSMixer at 120D, and Elastic Net at 5D.

The greater number of significant comparisons against Zero-Return suggests that the statistical conclusions are also sensitive to benchmark selection. However, the number of significant results remains limited relative to the full set of comparisons. These findings do not establish uniform predictive superiority across models or forecast horizons.

\paragraph{Benchmark sensitivity across volatility regimes.} Table~\ref{tab:benchmark_regimes} further examines the differences between the two benchmark specifications across low-, normal-, and high-volatility regimes. At shorter horizons, the average differences are relatively small. At longer horizons, however, the gaps become more pronounced, particularly under low- and normal-volatility conditions. At 120D, the difference reaches 30.79 percentage points in the low-volatility regime and 37.93 percentage points in the normal-volatility regime, compared with 16.28 percentage points in the high-volatility regime.

The regime-level results reinforce the finding that benchmark selection has a greater influence on measured performance at longer horizons. They also indicate that the magnitude of this effect varies across market conditions rather than remaining constant across volatility regimes.

Overall, the robustness analysis demonstrates that measured predictive performance is sensitive to the benchmark specification, especially at longer forecast horizons. Although several models achieve positive \(R^2_{\mathrm{OOS}}\) values and statistically significant improvements relative to Zero-Return, their performance is generally less favorable when assessed against HM. The primary conclusions of this study therefore continue to rely on the Historical Mean benchmark, while the Zero-Return results serve as a complementary robustness check.

% Table X: OOS R-squared under Historical Mean and Zero-Return benchmarks
\begin{table}[htbp]
\centering
\caption{Out-of-Sample $R^2$ under Alternative Benchmark Specifications}
\label{tab:benchmark_comparison}
\small
\setlength{\tabcolsep}{5pt}
\renewcommand{\arraystretch}{1.12}
\begin{tabular}{lrrrrr}
\hline
\textbf{Model} & \textbf{1D} & \textbf{5D} & \textbf{20D} & \textbf{60D} & \textbf{120D} \\
\hline
\multicolumn{6}{l}{\textit{Panel A: Historical Mean Benchmark (\%)}} \\
\hline
ARIMA          & -0.487 & -8.248 & -12.701 & -10.263 &  1.950 \\
ARIMAX         &  1.181 & -5.904 &  -4.610 &   2.518 & -0.349 \\
Ridge          &  1.498 &  1.786 &   1.657 &   0.906 & -4.947 \\
Elastic Net    &  0.127 &  0.609 &   0.000 &   0.000 &  0.000 \\
SVR            &  1.985 &  1.479 &   3.048 &   4.639 & -19.097 \\
Random Forest  &  1.771 & -1.297 &  -0.393 &   0.331 & -6.417 \\
XGBoost        &  1.478 &  0.109 &   1.374 &   0.130 & -0.141 \\
LSTM           & -0.651 & -0.505 &   1.475 &  -0.959 & -3.337 \\
Transformer    & -15.280 & -39.717 & -11.637 & -21.282 & -23.024 \\
TCN            & -14.156 & -64.910 & -123.428 & -87.232 & -142.548 \\
PatchTST       & -4.251 & -6.853 & -12.171 & -21.863 & -12.215 \\
TSMixer        & -2.719 & -5.472 &   3.795 &   4.982 & 17.535 \\
\hline
\multicolumn{6}{l}{\textit{Panel B: Zero-Return Benchmark (\%)}} \\
\hline
ARIMA          & -0.286 & -7.102 &  -7.414 &   6.408 & 27.033 \\
ARIMAX         &  1.379 & -4.782 &   0.298 &  17.257 & 25.322 \\
Ridge          &  1.695 &  2.826 &   6.270 &  15.888 & 21.900 \\
Elastic Net    &  0.327 &  1.662 &   4.691 &  15.119 & 25.582 \\
SVR            &  2.181 &  2.523 &   7.596 &  19.057 & 11.370 \\
Random Forest  &  1.968 & -0.224 &   4.316 &  15.400 & 20.806 \\
XGBoost        &  1.676 &  1.167 &   6.001 &  15.229 & 25.476 \\
LSTM           & -0.449 &  0.560 &   6.097 &  14.305 & 23.098 \\
Transformer    & -15.049 & -38.237 &  -6.400 &  -2.945 &  8.447 \\
TCN            & -13.927 & -63.163 & -112.947 & -58.924 & -80.500 \\
PatchTST       & -4.042 & -5.722 &  -6.909 &  -3.438 & 16.491 \\
TSMixer        & -2.513 & -4.355 &   8.308 &  19.348 & 38.631 \\
\hline
\end{tabular}
\par\vspace{4pt}\parbox{\linewidth}{\fontsize{9pt}{12pt}\selectfont\textit{Notes:} Values are expressed as percentages. Positive values indicate lower squared forecast errors than the corresponding benchmark, whereas negative values indicate higher errors. The Historical Mean benchmark is the primary reference; Zero-Return is included as an alternative benchmark specification.}
\end{table}
\FloatBarrier


% Table Y: Difference between Zero-Return and Historical Mean R-squared
\begin{table}[htbp]
\centering
\caption{Difference in Out-of-Sample $R^2$ between Zero-Return and Historical Mean Benchmarks}
\label{tab:benchmark_difference}
\small
\setlength{\tabcolsep}{6pt}
\renewcommand{\arraystretch}{1.12}
\begin{tabular}{lrrrrr}
\hline
\textbf{Model} & \textbf{1D} & \textbf{5D} & \textbf{20D} & \textbf{60D} & \textbf{120D} \\
\hline
ARIMA          & 0.201 & 1.146 &  5.287 & 16.671 & 25.083 \\
ARIMAX         & 0.198 & 1.122 &  4.907 & 14.738 & 25.671 \\
Ridge          & 0.197 & 1.040 &  4.613 & 14.982 & 26.847 \\
Elastic Net    & 0.200 & 1.053 &  4.691 & 15.119 & 25.582 \\
SVR            & 0.196 & 1.043 &  4.548 & 14.418 & 30.467 \\
Random Forest  & 0.197 & 1.073 &  4.709 & 15.069 & 27.223 \\
XGBoost        & 0.197 & 1.058 &  4.627 & 15.100 & 25.618 \\
LSTM           & 0.202 & 1.064 &  4.622 & 15.264 & 26.435 \\
Transformer    & 0.231 & 1.480 &  5.237 & 18.337 & 31.472 \\
TCN            & 0.229 & 1.746 & 10.481 & 28.308 & 62.048 \\
PatchTST       & 0.209 & 1.132 &  5.262 & 18.425 & 28.706 \\
TSMixer        & 0.206 & 1.117 &  4.513 & 14.366 & 21.096 \\
\hline
\textbf{Average} & \textbf{0.205} & \textbf{1.173} & \textbf{5.291} & \textbf{16.733} & \textbf{29.687} \\
\hline
\end{tabular}
\par\vspace{4pt}\parbox{\linewidth}{\fontsize{9pt}{12pt}\selectfont\textit{Notes:} Entries represent percentage-point differences calculated as $R^2_{\mathrm{OOS,Zero}}-R^2_{\mathrm{OOS,HM}}$. The overall average difference across all 60 model--horizon combinations is 10.618 percentage points.}
\end{table}

% Table Z: Summary of Diebold-Mariano test results
\begin{table}[htbp]
\centering
\caption{Summary of Diebold--Mariano Test Results across Benchmarks}
\label{tab:dm_benchmark_summary}
\small
\begin{tabular}{lrr}
\hline
\textbf{Benchmark} & \textbf{Significant comparisons} & \textbf{Share} \\
\hline
Historical Mean & 1/60 & 1.7\% \\
Zero-Return     & 8/60 & 13.3\% \\
\hline
\end{tabular}
\par\vspace{4pt}\parbox{\linewidth}{\fontsize{9pt}{12pt}\selectfont\textit{Notes:} One-sided Diebold--Mariano tests with the Harvey--Leybourne--Newbold correction. Significance is assessed at the reported 5\% level.}
\end{table}

% Table W: Benchmark comparison across volatility regimes
\begin{table}[htbp]
\centering
\caption{Benchmark Sensitivity across Volatility Regimes}
\label{tab:benchmark_regimes}
\small
\setlength{\tabcolsep}{7pt}
\renewcommand{\arraystretch}{1.15}
\begin{tabular}{llrrr}
\hline
\textbf{Horizon} & \textbf{Volatility Regime}
& \textbf{HM} & \textbf{Zero-Return} & \textbf{Difference} \\
& & \textbf{(\%)} & \textbf{(\%)} & \textbf{(pp)} \\
\hline
1D   & Low    & -4.04  & -3.98  &  0.07 \\
1D   & Normal & -4.41  & -4.07  &  0.34 \\
1D   & High   & -1.26  & -1.10  &  0.17 \\
\hline
5D   & Low    & -10.73 & -8.91  &  1.82 \\
5D   & Normal & -8.57  & -6.63  &  1.93 \\
5D   & High   & -11.81 & -11.13 &  0.68 \\
\hline
20D  & Low    & -8.01  &  0.03  &  8.04 \\
20D  & Normal & -3.50  &  7.71  & 11.20 \\
20D  & High   & -18.41 & -17.63 &  0.78 \\
\hline
60D  & Low    & -2.35  & 17.28  & 19.63 \\
60D  & Normal &  0.52  & 26.54  & 26.02 \\
60D  & High   & -20.94 & -16.67 &  4.28 \\
\hline
120D & Low    & -7.63  & 23.16  & 30.79 \\
120D & Normal & -4.11  & 33.82  & 37.93 \\
120D & High   & -28.43 & -12.15 & 16.28 \\
\hline
\end{tabular}
\par\vspace{4pt}\parbox{\linewidth}{\fontsize{9pt}{12pt}\selectfont\textit{Notes:} HM denotes the Historical Mean benchmark. Difference is calculated as the Zero-Return $R^2_{\mathrm{OOS}}$ minus the HM $R^2_{\mathrm{OOS}}$, in percentage points. Values are averaged across the 12 candidate models within each horizon--regime combination.}
\end{table}

\FloatBarrier

\section{Conclusion}\label{sec:conclusion}

This study examines out-of-sample (OOS) equity return predictability across five forecast horizons, three volatility regimes, and chronological evaluation periods using twelve econometric, machine learning, and deep learning models. Rather than identifying a universally superior model, it investigates how predictive performance varies across investment horizons, market volatility conditions, and evaluation periods.

First, descriptive results reveal non-monotonic performance across forecast horizons. SVR and Ridge perform relatively well at short and intermediate horizons, whereas TSMixer achieves the strongest relative performance at longer horizons, including an OOS $R^2$ of 17.53\% at 120 days. However, formal evidence remains limited: under the baseline HAC specification, only ARIMA retains significance after Holm correction. This result is robust across the examined HAC bandwidths but is not confirmed by the synchronized moving-block bootstrap. Thus, horizon-dependent differences are observable but are not statistically robust across forecasting models.

Second, predictive performance varies across volatility regimes and horizons, with no regime consistently favoring all models. The Normal regime contains the largest proportion of positive model--horizon OOS $R^2$ outcomes. Nevertheless, Giacomini--White tests identify significant cases exclusively involving underperformance relative to the Historical Mean (HM). The horizon--regime interaction test also identifies significant differences only for TCN after multiple-testing correction. These findings suggest model-specific sensitivity to volatility conditions rather than systematic predictive advantages under particular regimes. The robustness analyses using alternative 10-day and 60-day rolling volatility windows further support the main RQ2 conclusions, although the exact regime assignments and the statistical significance of individual horizon--regime interactions exhibit some sensitivity to the window definition.

Third, temporal stability tests identify significant variation for selected model--horizon combinations. Holdout block-level tests detect differences for TCN, LSTM, and PatchTST, while the Giacomini--Rossi fluctuation test identifies instability for Transformer, ARIMA, and TCN. However, validation block-level results do not remain significant after multiple-testing correction. In contrast, TSMixer shows comparatively consistent performance across the three holdout blocks at the 120-day horizon. Temporal instability therefore appears model- and horizon-dependent rather than universal.

Overall, the findings provide qualified evidence that equity return predictability varies across horizons, volatility regimes, and chronological periods. However, statistical support differs across these dimensions: horizon effects are sensitive to inference procedures, significant horizon--regime interactions are concentrated in one architecture, and temporal instability affects only selected models. The results therefore do not establish universal relationships between market conditions and OOS predictive performance.

This study contributes an integrated evaluation framework combining multi-horizon forecasting, volatility-conditioned assessment, and chronological stability analysis. Comparing twelve models against the HM benchmark also demonstrates that greater architectural complexity does not necessarily produce more reliable forecasts. Importantly, the study distinguishes descriptive improvements in accuracy from statistically supported predictive gains while accounting for multiple comparisons and serial dependence.

Several limitations remain. First, overlapping cumulative return targets induce serial dependence, making inference sensitive to covariance estimation and bootstrap specifications. The disagreement between the baseline H1 HAC test and synchronized moving-block bootstrap further highlights this concern. Second, Elastic Net produces forecasts identical to HM at certain longer horizons, creating degenerate loss differentials and reducing effective degrees of freedom in some tests. Third, the analysis uses a single benchmark and 20-day realized volatility classification without systematically examining alternative specifications. Significant conditional predictive ability results must also be interpreted carefully, as they may indicate underperformance rather than predictive improvement. Finally, the study covers only twelve model implementations and the S\&P~500, limiting generalizability to other architectures and financial markets.

Future research could evaluate alternative benchmarks, volatility measures, regime definitions, and dependence-robust inference methods. Controlled architectural comparisons, feature ablation, and additional market indices could further establish whether the observed patterns extend beyond the S\&P~500. Incorporating transaction costs, portfolio allocation, and risk-adjusted performance measures would also clarify whether statistically detectable forecasting differences translate into economically meaningful investment outcomes.

\bibliographystyle{acm}
\bibliography{references}

\end{document}
